% Beamer slide deck for Mathematical Physics
% "Mathematical Physics"
% Chapter 1: Complex Analysis
% Self-contained slide text; local figures live in figs.

\documentclass[11pt]{beamer}

\usepackage[utf8]{inputenc}
\usepackage[T1]{fontenc}
\usepackage{lmodern}
\usepackage{amsmath,amssymb}
\usepackage{graphicx}
\usepackage{etoolbox}

% Do not split aligned derivations between continuation slides.  Such splits
% create orphan equation lines and large blank regions in Beamer.

% All figure assets are required and shipped with this source bundle.
\graphicspath{{figs/}}
\newcommand{\safeincludegraphics}[2][]{\includegraphics[#1]{#2}}

\usetheme{Madrid}

\author[]{Compiled by Jake W. Liu}
\title[Complex Analysis]{Chapter 1: Complex Analysis}
\subtitle{Mathematical Physics}
\date{\today}

\setbeamertemplate{navigation symbols}{}
\setbeamertemplate{frametitle continuation}[from second]

% Per-section TOC frames removed to keep the deck near target length.

\usepackage{Kyushu}


% Layout controls follow the supplied local theme. The font size is unchanged; only the
% excessive display-math whitespace and side margins are reduced.
\setbeamersize{text margin left=0.65cm,text margin right=0.65cm}
\setlength{\emergencystretch}{1em}
\newcommand{\mpcompactspacing}{%
  \setlength{\abovedisplayskip}{3pt plus 1pt minus 1pt}%
  \setlength{\belowdisplayskip}{3pt plus 1pt minus 1pt}%
  \setlength{\abovedisplayshortskip}{2pt plus 1pt minus 1pt}%
  \setlength{\belowdisplayshortskip}{2pt plus 1pt minus 1pt}%
  \setlength{\jot}{2pt}%
}
\apptocmd{\normalsize}{\mpcompactspacing}{}{}
\apptocmd{\small}{\mpcompactspacing}{}{}
\apptocmd{\footnotesize}{\mpcompactspacing}{}{}
\apptocmd{\scriptsize}{\mpcompactspacing}{}{}

% Neutralize \index so source markup (if any) is harmless.
\makeatletter
\@ifundefined{index}{\newcommand{\index}[1]{}}{\renewcommand{\index}[1]{}}
\makeatother

\begin{document}

\newcounter{sfchapter}
\setcounter{sfchapter}{1}
\renewcommand{\thesection}{\arabic{sfchapter}.\arabic{section}}
\numberwithin{equation}{section}

\begin{frame}[c]\titlepage
\end{frame}

\begin{frame}{Outline}
\tableofcontents
\end{frame}

\begin{frame}{Where we are going}
\footnotesize
To analyze the special functions of physics we first build the tools of
\emph{complex analysis}.
\begin{itemize}
  \item \textbf{Taylor and Laurent series} represent a function near a point, even
        near a singularity.
  \item \textbf{Analytic functions} (complex differentiable) obey the
        \emph{Cauchy--Riemann equations} and admit power-series expansions.
  \item \textbf{Cauchy's integral theorem} and the \textbf{residue theorem} turn contour
        integrals into a sum of local residues---the tool used later for Gamma/Beta identities
        and real integrals.
  \item A short catalog of \textbf{contour templates} (semicircle, keyhole, Jordan) shows
        how orientation and vanishing-arc estimates convert complex residues into real
        answers.
  \item The \textbf{method of steepest descent} extracts large-parameter asymptotics of
        contour integrals. It gives large-argument Bessel expansions (Ch.~5); its real
        special cases (Laplace method, stationary phase) give Stirling (Ch.~2) and Airy (Ch.~3).
\end{itemize}
\end{frame}

% ======================================================================
\section{Taylor series representation}
% ======================================================================

\begin{frame}[allowframebreaks=1.00]{Matching a polynomial to a function}
\footnotesize

A Taylor polynomial replaces a function near one point by a polynomial with the
same value, slope, curvature, and higher derivatives there.

Start at $x=0$ and write the first few terms rather than a general sum:
\[
 p(x)=a_0+a_1x+a_2x^2+a_3x^3+\cdots+a_Nx^N\qquad(N\ge3).
\]
Differentiate and then set $x=0$:
\begin{align*}
 p(0)&=a_0,\\
 p'(x)&=a_1+2a_2x+3a_3x^2+\cdots,
       &p'(0)&=a_1,\\
 p''(x)&=2a_2+6a_3x+\cdots,
       &p''(0)&=2!a_2,\\
 p^{(3)}(x)&=6a_3+\cdots,
       &p^{(3)}(0)&=3!a_3.
\end{align*}
To match $f$, choose $a_0=f(0)$, $a_1=f'(0)$, $a_2=f''(0)/2!$, and so on.
The pattern is
\begin{equation}
 a_k=\frac{f^{(k)}(0)}{k!}.
 \label{eq:taylor-coef}
\end{equation}
The factorial appears because differentiating $x^k$ exactly $k$ times gives $k!$.


At a different center $x_0$, use powers of the distance $x-x_0$:
\begin{equation}
 T_N(x;x_0)=\sum_{k=0}^{N}\frac{f^{(k)}(x_0)}{k!}(x-x_0)^k.
 \label{eq:taylorx0}
\end{equation}


\textbf{When does the infinite series equal \(f\)?} Let \(N\to\infty\). The error of the
approximation is the remainder \(R_N(x)=f(x)-T_N(x;x_0)\). If \(R_N(x)\to0\), then \(f\) equals its
Taylor series at \(x\):
\begin{equation}
 f(x)=\sum_{k=0}^{\infty}\frac{f^{(k)}(x_0)}{k!}(x-x_0)^k.
 \label{eq:taylor0}
\end{equation}
\par\framebreak\par
\textbf{Warning: smooth is not enough.} Here is a function with derivatives of every order whose
Taylor series is wrong. Let
\[
 f(x)=e^{-1/x^2}\ (x\ne0),\qquad f(0)=0 .
\]
\begin{itemize}
  \item \emph{Derivatives away from \(0\).} By the chain rule, \(f'(x)=\frac{2}{x^3}e^{-1/x^2}\).
        Differentiating again and again always gives (a polynomial in \(1/x\))\(\times e^{-1/x^2}\).
  \item \emph{Each such term tends to \(0\) as \(x\to0\).} Put \(y=1/x^2\to\infty\). A term
        \(|x|^{-m}e^{-1/x^2}\) becomes \(y^{m/2}e^{-y}\), and \(e^{-y}\) beats every power of \(y\).
  \item \emph{Derivatives at \(0\).} By definition \(f^{(k)}(0)=\lim_{h\to0}f^{(k-1)}(h)/h\). This is
        again (a polynomial in \(1/h\))\(\times e^{-1/h^2}\), so it is \(0\). Hence
        \(f^{(k)}(0)=0\) for every \(k\).
\end{itemize}
So the Taylor series of \(f\) at \(0\) is \(0+0\cdot x+0\cdot x^2+\cdots=0\). It converges everywhere,
but it equals \(f(x)\) at no \(x\neq0\), since \(f(x)>0\) there. Such an \(f\) is called \emph{not
analytic} at \(0\). Section~1.5 shows that complex-differentiable functions never behave like
this.
\par\framebreak\par


\textbf{Exponential.} Every derivative of $e^x$ equals $e^x$, so every derivative at
zero equals $1$. Therefore
\begin{equation}
 e^x=\sum_{k=0}^{\infty}\frac{x^k}{k!}
 =1+x+\frac{x^2}{2!}+\frac{x^3}{3!}+\cdots.
 \label{eq:exp-series}
\end{equation}
The denominator $k!$ grows faster than any power $x^k$, so the terms shrink rapidly
and the series converges for every $x$, where it represents $e^x$.
For example, at $x=0.1$ the first three terms give
$e^{0.1}\approx1+0.1+0.005=1.105$.


\textbf{Cosine and sine.} List the derivative values at zero and divide by the
corresponding factorial:
\[
\begin{array}{c|rrrrrr}
 k&0&1&2&3&4&5\\ \hline
 (\cos)^{(k)}(0)&1&0&-1&0&1&0\\
 (\sin)^{(k)}(0)&0&1&0&-1&0&1
\end{array}
\]
Only even powers survive for cosine; only odd powers survive for sine:
\begin{equation}
 \cos x=\sum_{k=0}^{\infty}\frac{(-1)^k x^{2k}}{(2k)!}
 =1-\frac{x^2}{2!}+\frac{x^4}{4!}-\cdots,
 \label{eq:cos-series}
\end{equation}
\begin{equation}
 \sin x=\sum_{k=0}^{\infty}\frac{(-1)^k x^{2k+1}}{(2k+1)!}
 =x-\frac{x^3}{3!}+\frac{x^5}{5!}-\cdots.
 \label{eq:sin-series}
\end{equation}
Every derivative of $\cos x$ and $\sin x$ has magnitude at most $1$, so the same
factorial dominance makes both series converge for every $x$.
All three series also converge for every complex argument (the same $k!$ comparison works
with $|z|$ in place of $|x|$). We use them to extend
$e^x$, $\cos x$, and $\sin x$ to complex numbers.
\end{frame}

% ======================================================================
\section{Complex variables}
% ======================================================================

\begin{frame}[allowframebreaks=1.00]{Imaginary unit, Cartesian and polar form}
\footnotesize

Define the imaginary unit $i$ by $i^2=-1$; the other square root of $-1$ is $-i$.
A complex number has Cartesian form
\[
  z=x+iy,\qquad x=\operatorname{Re}z,\quad y=\operatorname{Im}z,
\]
where $x,y\in\mathbb R$; $\operatorname{Re}z$ and $\operatorname{Im}z$ denote the real
and imaginary parts of $z$.  The powers of $i$ repeat with period four:
\[
  i^{4q}=1,\qquad i^{4q+1}=i,\qquad
  i^{4q+2}=-1,\qquad i^{4q+3}=-i.
\]
Substitute $i\theta$ into the exponential series \eqref{eq:exp-series} and separate even
and odd powers:
\begin{align*}
 e^{i\theta}
  &=\sum_{n=0}^{\infty}\frac{(i\theta)^n}{n!}\\
  &=\sum_{k=0}^{\infty}\frac{i^{2k}\theta^{2k}}{(2k)!}
    +\sum_{k=0}^{\infty}\frac{i^{2k+1}\theta^{2k+1}}{(2k+1)!}\\
 \intertext{and since $i^{2k}=(i^2)^k=(-1)^k$ and $i^{2k+1}=i\,i^{2k}=i(-1)^k$,}
  &=\sum_{k=0}^{\infty}\frac{(-1)^k\theta^{2k}}{(2k)!}
    +i\sum_{k=0}^{\infty}\frac{(-1)^k\theta^{2k+1}}{(2k+1)!}.
\end{align*}
Comparison with \eqref{eq:cos-series}--\eqref{eq:sin-series} gives Euler's formula
\begin{equation}
  e^{i\theta}=\cos\theta+i\sin\theta.
  \label{eq:euler}
\end{equation}

\textbf{The addition law \(e^{w_1+w_2}=e^{w_1}e^{w_2}\) for complex numbers.} For real numbers we
know it; for complex numbers we must check it from the series, because \(e^w\) is now
\emph{defined} by the series.
\begin{itemize}
  \item \emph{Left side.} \(e^{w_1+w_2}=\sum_n (w_1+w_2)^n/n!\). By the binomial theorem,
        \((w_1+w_2)^n=\sum_{k=0}^n\frac{n!}{k!\,(n-k)!}w_1^kw_2^{n-k}\), so
        \[
         \frac{(w_1+w_2)^n}{n!}=\sum_{k=0}^{n}\frac{w_1^k}{k!}\,\frac{w_2^{\,n-k}}{(n-k)!}.
        \]
        Summing over \(n\) gives every term \(\frac{w_1^k}{k!}\frac{w_2^j}{j!}\) with \(k,j\ge0\) exactly
        once (the one with \(n=k+j\)).
  \item \emph{Right side.} Multiplying out
        \(\bigl(\sum_k w_1^k/k!\bigr)\bigl(\sum_j w_2^j/j!\bigr)\) gives the same terms
        \(\frac{w_1^k}{k!}\frac{w_2^j}{j!}\), each once.
\end{itemize}
So the two sides are equal. (Rearranging the terms is allowed because the series converge
absolutely.) We use this to split \(e^{u+iv}\), which the next frame needs. Take \(w_1=u\),
\(w_2=iv\), and use \eqref{eq:euler} for \(e^{iv}\):
\begin{equation}
 e^{u+iv}=e^u(\cos v+i\sin v),\qquad u,v\in\mathbb R.
 \label{eq:exp-uv}
\end{equation}
So \(|e^{u+iv}|=e^u\) (since \(|\cos v+i\sin v|=\sqrt{\cos^2v+\sin^2v}=1\)), and \(v\) is the angle.
\par\framebreak\par
\textbf{Polar form.} For \(z\ne0\), set
\[
  r=|z|=\sqrt{x^2+y^2}>0.
\]
Choose any angle $\theta$ satisfying
$x=r\cos\theta$ and $y=r\sin\theta$.  Then
\begin{equation}
  z=r(\cos\theta+i\sin\theta)=re^{i\theta}.
  \label{eq:polar}
\end{equation}
The angle is not unique: if $\theta$ is one argument, then every argument is
$\theta+2\pi k$, $k\in\mathbb Z$.  The principal argument is commonly chosen as
$\operatorname{Arg}z\in(-\pi,\pi]$.
The number $z=0$ has modulus $0$ but no defined argument.

\textbf{Products in polar form.} If \(z_1=r_1e^{i\theta_1}\) and \(z_2=r_2e^{i\theta_2}\), the
addition law gives \(e^{i\theta_1}e^{i\theta_2}=e^{i(\theta_1+\theta_2)}\), so
\[
  z_1z_2=r_1r_2\,e^{i\theta_1}e^{i\theta_2}=r_1r_2\,e^{i(\theta_1+\theta_2)} .
\]
In words: to multiply complex numbers, multiply their moduli and add their angles.
\par\framebreak\par
\textbf{Complex conjugate.} The complex conjugate is
\[
  \overline z=x-iy=re^{-i\theta},
  \qquad
  z\overline z=(x+iy)(x-iy)=x^2+y^2=r^2=|z|^2,
\]
Adding \(z+\overline z=2x\) and subtracting \(z-\overline z=2iy\) isolates the real and imaginary
parts:
\begin{equation}
  \operatorname{Re}z=\frac{z+\overline z}{2},\qquad
  \operatorname{Im}z=\frac{z-\overline z}{2i}.
  \label{eq:re-im-conj}
\end{equation}
\textbf{Cosine and sine from exponentials.} Replace \(\theta\) by \(-\theta\) in Euler's formula
\eqref{eq:euler}. Since \(\cos(-\theta)=\cos\theta\) and \(\sin(-\theta)=-\sin\theta\),
\[
 e^{-i\theta}=\cos\theta-i\sin\theta .
\]
Add this to \(e^{i\theta}=\cos\theta+i\sin\theta\), and subtract it:
\begin{align*}
  e^{i\theta}+e^{-i\theta}&=2\cos\theta,\\
  e^{i\theta}-e^{-i\theta}&=2i\sin\theta,
\end{align*}
and therefore
\begin{equation}
  \cos\theta=\frac{e^{i\theta}+e^{-i\theta}}{2},\qquad
  \sin\theta=\frac{e^{i\theta}-e^{-i\theta}}{2i}.
  \label{eq:cos-sin-exp}
\end{equation}
\end{frame}



\begin{frame}[allowframebreaks=1.00]{Multivalued functions and branch cuts}
\footnotesize

\textbf{Why we need this.} Powers \(z^a\) and the logarithm \(\log z\) appear in many integrals later.
(In the keyhole contour of Section~1.8 the whole answer comes from the jump of \(z^a\) across a
cut.) These functions can take several values at one point, and we must learn how to choose one.

\textbf{A familiar example.} The equation \(x^2=y\) (with \(y>0\)) has two solutions,
\(x=\pm\sqrt y\). The real square root picks one of them, \(+\sqrt y\), in a consistent way. A
\emph{branch} of a many-valued function is such a consistent choice on a chosen region.

\textbf{The complex logarithm.} We want to solve \(e^w=z\) for \(w\), given \(z\ne0\). Write
\(z=re^{i\theta}\) with \(r>0\), and \(w=u+iv\) with \(u,v\) real. By \eqref{eq:exp-uv},
\[
 e^w=e^{u+iv}=e^u(\cos v+i\sin v).
\]
We need this to equal \(r(\cos\theta+i\sin\theta)\). Compare moduli and angles separately:
\begin{itemize}
  \item moduli: \(e^u=r\), so \(u=\log r\) (the ordinary real logarithm, since \(r>0\));
  \item angles: \(\cos v=\cos\theta\) and \(\sin v=\sin\theta\), so \(v=\theta+2\pi k\) for some
        integer \(k\). Every \(k\) works, because \(\cos\) and \(\sin\) repeat with period \(2\pi\).
\end{itemize}
So \(e^w=z\) has infinitely many solutions, one for each integer \(k\):
\begin{equation}
  \log z=\log|z|+i(\theta+2\pi k),\qquad k\in\mathbb Z.
  \label{eq:complex-log-values}
\end{equation}
\textbf{Choosing one value: the principal branch.} Restrict the angle to one interval of length
\(2\pi\), usually \(-\pi<\operatorname{Arg}z<\pi\), and define
\[
  \operatorname{Log}z=\log|z|+i\operatorname{Arg}z .
\]
This is a single-valued, analytic function on the plane with the half-line \((-\infty,0]\) removed.
\par\framebreak\par
\textbf{Why something must be removed: the branch point.} Walk once counterclockwise around
\(0\) on the circle \(z=re^{i\theta}\), letting \(\theta\) increase continuously from \(0\) to \(2\pi\).
Follow \(\log z=\log r+i\theta\) continuously along the way: it starts at \(\log r\) and ends at
\(\log r+2\pi i\). We are back at the same point \(z\), but the logarithm did not come back to the
same value. A point with this property is a \emph{branch point}; here it is \(z=0\).

\textbf{The branch cut.} To get a single-valued function, we forbid such loops by removing a curve
from the branch point to infinity, a \emph{branch cut}. For \(\operatorname{Log}\) the cut is the
negative real axis. Just above and just below the cut, at \(-r\) (\(r>0\)),
\[
  \operatorname{Log}(-r+i0)=\log r+i\pi,\qquad
  \operatorname{Log}(-r-i0)=\log r-i\pi .
\]
So \(\operatorname{Log}\) jumps by \(2\pi i\) across the cut.

\textbf{Roots work the same way.} Define \(z^{1/2}=e^{\frac12\log z}\) and insert all the values
\eqref{eq:complex-log-values}:
\[
 z^{1/2}=e^{\frac12\log r}\,e^{i(\theta+2\pi k)/2}
       =\sqrt r\,e^{i\theta/2}e^{i\pi k}
       =(-1)^k\sqrt r\,e^{i\theta/2},
\]
since \(e^{i\pi k}=\cos\pi k+i\sin\pi k=(-1)^k\) by \eqref{eq:euler}. So \(z^{1/2}\) has exactly two
values, \(\pm\sqrt r\,e^{i\theta/2}\) (even and odd \(k\)), just like \(\pm\sqrt y\). A branch cut lets us
pick one of them consistently. In general \(z^a=e^{a\log z}\), and the same cut is used.
\end{frame}



% ======================================================================
\section{Derivatives}
% ======================================================================

\begin{frame}[allowframebreaks=1.00]{The complex derivative and Cauchy--Riemann}
\footnotesize

\textbf{Definition.} The complex derivative of \(f\) at \(z_0\) is
\begin{equation}
 f'(z_0)=\lim_{h\to0}\frac{f(z_0+h)-f(z_0)}{h}.
 \label{eq:cderiv}
\end{equation}
It looks like the real definition, but there is one big difference. Here \(h\) is a complex
number, so it can approach \(0\) from \emph{any direction} in the plane (along the real axis,
along the imaginary axis, along a spiral, \dots). The derivative exists only if every direction
gives the same limit. This is a strong requirement, and it is what makes complex analysis so
powerful.

\textbf{Picture.} For small \(h\), \(f(z_0+h)\approx f(z_0)+f'(z_0)\,h\). So near \(z_0\), \(f\) multiplies
every small displacement \(h\) by the same complex number \(f'(z_0)\). Multiplying by a complex
number means rotating by its angle \(\arg f'(z_0)\) and stretching by its size \(|f'(z_0)|\) (polar
form). So a differentiable \(f\) looks, up close, like a rotation plus a uniform stretch.

\textbf{Names.} \(f\) is \emph{holomorphic} (or \emph{analytic}) on an open set if \(f'\) exists at
every point of the set. It is \emph{entire} if it is holomorphic on the whole plane \(\mathbb C\).
\par\framebreak\par
\textbf{Goal: a practical test.} Write \(z=x+iy\) and split \(f\) into real and imaginary parts,
\(f(z)=u(x,y)+iv(x,y)\). Subscripts denote partial derivatives (\(u_x=\partial u/\partial x\),
etc.). We compute the limit \eqref{eq:cderiv} in two directions and demand that they agree.

\textbf{Direction 1: \(h\) real}, \(h=\Delta x\). Only \(x\) changes:
\[
 \frac{f(z+h)-f(z)}{h}
 =\frac{u(x+\Delta x,y)-u(x,y)}{\Delta x}
  +i\,\frac{v(x+\Delta x,y)-v(x,y)}{\Delta x}.
\]
As \(\Delta x\to0\) the two quotients become partial derivatives in \(x\): the limit is
\(u_x+iv_x\).

\textbf{Direction 2: \(h\) imaginary}, \(h=i\Delta y\). Only \(y\) changes, and the denominator is
\(i\Delta y\):
\[
 \frac{f(z+i\Delta y)-f(z)}{i\,\Delta y}
 =\frac1i\left[\frac{u(x,y+\Delta y)-u(x,y)}{\Delta y}
   +i\,\frac{v(x,y+\Delta y)-v(x,y)}{\Delta y}\right].
\]
As \(\Delta y\to0\) the quotients become partial derivatives in \(y\): the limit is
\(\frac1i\bigl(u_y+iv_y\bigr)\).
Simplify with \(1/i=-i\): \(\;-i(u_y+iv_y)=-iu_y-i^2v_y=v_y-iu_y\).
\par\framebreak\par
\textbf{Demand agreement.} If \(f'(z)\) exists, both directions give it:
\[
 u_x+iv_x=v_y-iu_y .
\]
Two complex numbers are equal when their real parts and their imaginary parts are equal. This
gives the \emph{Cauchy--Riemann equations}:
\begin{equation}
  u_x=v_y,\qquad u_y=-v_x.
  \label{eq:CR}
\end{equation}
So they are \emph{necessary}: every complex-differentiable function satisfies them, and then
\(f'=u_x+iv_x\).

\textbf{They are also sufficient}, if the four partial derivatives are continuous. Reason: for a
small step \(h=\Delta x+i\Delta y\), the changes of \(u\) and \(v\) are (to first order)
\[
 \Delta u\approx u_x\Delta x+u_y\Delta y,\qquad \Delta v\approx v_x\Delta x+v_y\Delta y .
\]
Replace \(u_y\) by \(-v_x\) and \(v_y\) by \(u_x\):
\[
 \Delta f=\Delta u+i\Delta v\approx u_x\Delta x-v_x\Delta y+iv_x\Delta x+iu_x\Delta y
 =(u_x+iv_x)(\Delta x+i\Delta y).
\]
(Multiply out the right side to check.) So \(\Delta f\approx(u_x+iv_x)\,h\) for every direction of
\(h\): the derivative exists and equals \(u_x+iv_x\).
\par\framebreak\par
\textbf{How to use the test.} Find \(u\) and \(v\), compute the four partial derivatives, and check
\eqref{eq:CR}.

\textbf{Example: \(f(z)=z^2\).} Expand: \(z^2=(x+iy)^2=x^2+2ixy+i^2y^2=(x^2-y^2)+i(2xy)\). So
\[
 u=x^2-y^2,\quad v=2xy:\qquad
 u_x=2x,\quad v_y=2x,\qquad u_y=-2y,\quad v_x=2y .
\]
Both equations hold everywhere, so \(z^2\) is entire, and
\(f'(z)=u_x+iv_x=2x+2iy=2z\), as expected.

\textbf{Counterexample: \(f(z)=\overline z=x-iy\).} Here \(u=x\), \(v=-y\), so \(u_x=1\) but
\(v_y=-1\). The first equation fails at every point: \(\overline z\) is nowhere differentiable. You can
also see it directly: the quotient \((\overline{z+h}-\overline z)/h=\overline h/h\) equals \(1\) for real \(h\)
and \(-1\) for imaginary \(h\), so the limit depends on the direction.
\end{frame}

% ======================================================================
\section{Path integrals}
% ======================================================================

\begin{frame}[allowframebreaks=1.00]{Contour integrals and a key circle integral}
\footnotesize

\textbf{What a contour integral is.} It is an ordinary integral along a chosen path in the
plane. Describe the path by a parameter: \(z=\gamma(t)\) for \(a\le t\le b\) (with \(\gamma\) smooth, or
smooth in pieces). Along the path \(dz=\gamma'(t)\,dt\), so
\begin{equation}
 \int_\gamma f(z)\,dz=\int_a^b f(\gamma(t))\,\gamma'(t)\,dt.
 \label{eq:pathint}
\end{equation}
The right side is a normal integral in \(t\) (of a complex-valued function: integrate the real and
imaginary parts separately).

\textbf{Direction matters.} Running the same path backwards changes the sign:
\(\int_{-\gamma}f\,dz=-\int_\gamma f\,dz\). The \emph{positive} direction around a region keeps the
region on your left: counterclockwise around the outside, clockwise around a hole.
\par\framebreak\par
\textbf{The key example: integer powers on a circle.} This simple computation is behind almost
everything later (Laurent coefficients, residues). Let \(n\) be an integer and integrate \(z^n\)
counterclockwise around the circle \(|z|=R\).
\begin{itemize}
  \item \emph{Parametrize:} \(z=Re^{i\theta}\), \(0\le\theta\le2\pi\), so \(dz=iRe^{i\theta}\,d\theta\) and
        \(z^n=R^ne^{in\theta}\).
  \item \emph{Substitute} into \eqref{eq:pathint} and combine the exponentials:
        \[
         \oint_{|z|=R}z^n\,dz
         =\int_0^{2\pi}R^ne^{in\theta}\cdot iRe^{i\theta}\,d\theta
         =iR^{n+1}\int_0^{2\pi}e^{i(n+1)\theta}\,d\theta.
        \]
\end{itemize}
\textbf{Case \(n=-1\).} Then \(e^{i(n+1)\theta}=e^0=1\) and \(R^{n+1}=1\):
\(\displaystyle\oint_{|z|=R}\frac{dz}{z}=i\int_0^{2\pi}d\theta=2\pi i\).

\textbf{Case \(n\ne-1\).} Integrate the exponential:
\[
 iR^{n+1}\int_0^{2\pi}e^{i(n+1)\theta}\,d\theta
 =iR^{n+1}\left[\frac{e^{i(n+1)\theta}}{i(n+1)}\right]_{0}^{2\pi}
 =\frac{R^{n+1}}{n+1}\bigl(e^{2\pi i(n+1)}-1\bigr)=0,
\]
because \(n+1\) is an integer, so \(e^{2\pi i(n+1)}=\cos(2\pi(n+1))+i\sin(2\pi(n+1))=1\).
\par\framebreak\par
\textbf{Result, for a circle around any center \(z_0\).} Substitute \(z=z_0+w\) (so \(dz=dw\), and the
circle \(|z-z_0|=R\) becomes \(|w|=R\)). Writing \(\delta_{n,-1}\) for \(1\) if \(n=-1\) and \(0\) otherwise,
\begin{equation}
 \oint_{|z-z_0|=R}(z-z_0)^n\,dz=2\pi i\,\delta_{n,-1}.
 \label{eq:circ-power}
\end{equation}
Notice that the answer does not depend on the radius \(R\).

\textbf{Why only \(n=-1\) survives.} For \(n\ne-1\), \(z^n\) has the antiderivative
\(z^{n+1}/(n+1)\), which takes the same value at the start and the end of a closed loop; so the
integral is \(0\). The antiderivative of \(1/z\) is \(\log z\) (check: differentiating \(e^{\log z}=z\)
gives \(e^{\log z}\cdot(\log z)'=1\), so \((\log z)'=1/z\)). But \(\log z\) gains \(2\pi i\) once around the
origin (Section~1.2), and that is exactly the value of the integral.
\end{frame}


\begin{frame}[allowframebreaks=1.00]{Cauchy's integral theorem and formula}
\footnotesize

\textbf{The question.} The circle computation showed that \(z^n\) with \(n\ge0\) integrates to zero
around a closed loop. Is that true for every holomorphic function? Cauchy's theorem says yes,
provided the loop does not surround any point where \(f\) fails to be holomorphic.

A connected open region is \emph{simply connected} if it has no holes (every closed curve in it
can be shrunk to a point without leaving the region).

\begin{block}{Cauchy's integral theorem}
Let $f$ be holomorphic on a simply connected region $D$, and let $\gamma$ be a
closed piecewise smooth contour in $D$.  Then
\[
  \oint_\gamma f(z)\,dz=0.
\]
\end{block}
\par\framebreak\par
\textbf{Why is it zero?} Let \(\gamma\) be the boundary of a region \(E\), traversed counterclockwise.
\begin{itemize}
  \item \emph{Step 1: split into real line integrals.} With \(f=u+iv\) and \(dz=dx+i\,dy\),
        \[
          f\,dz=(u+iv)(dx+i\,dy)=(u\,dx-v\,dy)+i\,(v\,dx+u\,dy).
        \]
  \item \emph{Step 2: Green's theorem}, \(\oint_{\partial E}(P\,dx+Q\,dy)=\iint_E(Q_x-P_y)\,dA\), applied
        to each part:
        \begin{itemize}
          \item real part, \(P=u\), \(Q=-v\): \(\;Q_x-P_y=-v_x-u_y\);
          \item imaginary part, \(P=v\), \(Q=u\): \(\;Q_x-P_y=u_x-v_y\).
        \end{itemize}
  \item \emph{Step 3: Cauchy--Riemann.} By \eqref{eq:CR}, \(v_x=-u_y\) and \(v_y=u_x\). So
        \(-v_x-u_y=u_y-u_y=0\) and \(u_x-v_y=u_x-u_x=0\).
\end{itemize}
Putting it together:
\[
 \oint_\gamma f\,dz
 =\iint_E(-v_x-u_y)\,dA+i\iint_E(u_x-v_y)\,dA
 =\iint_E0\,dA+i\iint_E0\,dA=0.
\]
(Green's theorem needs continuous partial derivatives; Goursat later showed that being
holomorphic is enough.)
\par\framebreak\par
\textbf{What happens if the loop surrounds a bad point?} Then the integral need not be zero; the
circle integral \(\oint dz/z=2\pi i\) is an example. The next result turns this into a tool.

\begin{block}{Cauchy's integral formula}
Let $\gamma$ be a positively oriented simple closed contour, let $f$ be holomorphic on an
open set containing $\gamma$ and its interior, and let $z_0$ lie inside $\gamma$.  Then
\begin{equation}
  f(z_0)=\frac{1}{2\pi i}\oint_\gamma\frac{f(z)}{z-z_0}\,dz.
  \label{eq:CIF}
\end{equation}
\end{block}
\textbf{In words:} the values of \(f\) on a loop determine its value at every point inside.
\par\framebreak\par
\begin{columns}[T]
\begin{column}{0.56\linewidth}
\textbf{Proof, step 1: shrink the loop.} The integrand \(f(z)/(z-z_0)\) is holomorphic except at
\(z_0\). Draw a small counterclockwise circle \(C_\varepsilon\) around \(z_0\) and join it to \(\gamma\) by a
slit \(s\) (right). The path \(\gamma+s-C_\varepsilon-s\) surrounds only the shaded region, where the
integrand is holomorphic, so its integral is \(0\) by Cauchy's theorem. The two sides of the slit
cancel. Hence
\[
 \oint_\gamma\frac{f(z)}{z-z_0}\,dz=\oint_{C_\varepsilon}\frac{f(z)}{z-z_0}\,dz .
\]
\end{column}
\begin{column}{0.4\linewidth}
\centering
\safeincludegraphics[width=\linewidth,height=0.3\textheight,keepaspectratio]{figs/slit-contour.pdf}
\end{column}
\end{columns}

\textbf{Step 2: compute on the small circle.} Parametrize \(z=z_0+\varepsilon e^{i\theta}\),
\(0\le\theta\le2\pi\). Then \(z-z_0=\varepsilon e^{i\theta}\) and \(dz=i\varepsilon e^{i\theta}\,d\theta\), so
\(\frac{dz}{z-z_0}=i\,d\theta\), and
\[
 \oint_{C_\varepsilon}\frac{f(z)}{z-z_0}\,dz
 =i\int_0^{2\pi}f(z_0+\varepsilon e^{i\theta})\,d\theta.
\]
\textbf{Step 3: let \(\varepsilon\to0\).} By continuity, \(f(z_0+\varepsilon e^{i\theta})\to f(z_0)\), so the
right side tends to \(i\cdot2\pi\cdot f(z_0)\). The left side does not depend on \(\varepsilon\) (Step~1).
So \(\oint_\gamma f(z)\,dz/(z-z_0)=2\pi i\,f(z_0)\), which is \eqref{eq:CIF}.
\par\framebreak\par
\textbf{A physical reading: the mean-value property.} Since the circle integral is the same for
every \(\varepsilon\), Step~2 and the formula give, for every such circle,
\[
 f(z_0)=\frac1{2\pi}\int_0^{2\pi}f(z_0+\varepsilon e^{i\theta})\,d\theta :
\]
the value at the center equals the average over the circle. The same is true of an
electrostatic potential in a region without charges.
\end{frame}



% ======================================================================
\section{Taylor and Laurent series}
% ======================================================================

\begin{frame}[allowframebreaks=1.00]{Taylor series from contour integration}
\footnotesize

\textbf{Goal.} Show that a holomorphic function equals its Taylor series in every disk where it
is holomorphic. A Taylor series needs all the derivatives of \(f\); Cauchy's formula supplies them.

\textbf{Step 1: derivatives from Cauchy's formula.} In \eqref{eq:CIF}, \(z_0\) appears only in
\(1/(z-z_0)\). Differentiate with respect to \(z_0\) (the contour stays fixed):
\[
 \frac{\partial}{\partial z_0}\frac1{z-z_0}=\frac1{(z-z_0)^2},\qquad
 \frac{\partial^2}{\partial z_0^2}\frac1{z-z_0}=\frac{2}{(z-z_0)^3},\qquad
 \frac{\partial^3}{\partial z_0^3}\frac1{z-z_0}=\frac{3!}{(z-z_0)^4},
\]
and in general \(\partial^n_{z_0}(z-z_0)^{-1}=n!\,(z-z_0)^{-(n+1)}\) (each derivative brings down the next
integer).
So, for \(n=0,1,2,\ldots\),
\begin{equation}
 f^{(n)}(z_0)=\frac{n!}{2\pi i}
  \oint_\gamma\frac{f(z)}{(z-z_0)^{n+1}}\,dz.
 \label{eq:CIF-deriv}
\end{equation}
The right side exists for every \(n\). So a function that is complex-differentiable once is
automatically differentiable infinitely many times. (Nothing like this holds for real functions.)
\par\framebreak\par
\begin{block}{Taylor series theorem}
If $f$ is holomorphic on the disk $|z-z_0|<R$, then for every $|z-z_0|<R$,
\[
 f(z)=\sum_{n=0}^{\infty}\frac{f^{(n)}(z_0)}{n!}(z-z_0)^n.
\]
\end{block}
\textbf{Step 2: Cauchy's formula on a circle.} Fix \(z\) in the disk. Choose a radius \(\rho\) with
\(|z-z_0|<\rho<R\), and let \(C_\rho\) be the circle \(|\xi-z_0|=\rho\). Cauchy's formula \eqref{eq:CIF},
written with integration variable \(\xi\) and interior point \(z\), gives
\[
 f(z)=\frac{1}{2\pi i}\oint_{C_\rho}\frac{f(\xi)}{\xi-z}\,d\xi.
\]
\par\framebreak\par
\textbf{Step 3: expand the kernel \(1/(\xi-z)\) in powers of \(z-z_0\).} Write
\(\xi-z=(\xi-z_0)-(z-z_0)\) and factor out \(\xi-z_0\):
\[
 \frac1{\xi-z}=\frac1{\xi-z_0}\cdot\frac1{1-w},\qquad w=\frac{z-z_0}{\xi-z_0}.
\]
On the circle \(|w|=|z-z_0|/\rho<1\), so the geometric series \(\frac1{1-w}=1+w+w^2+\cdots\) converges:
\[
 \frac1{\xi-z}=\frac1{\xi-z_0}\sum_{n=0}^{\infty}\Bigl(\frac{z-z_0}{\xi-z_0}\Bigr)^n
 =\sum_{n=0}^{\infty}\frac{(z-z_0)^n}{(\xi-z_0)^{n+1}}.
\]
\par\framebreak\par
\textbf{Step 4: integrate term by term.} Insert the series into Step~2. The factor \((z-z_0)^n\)
does not depend on \(\xi\), so it comes out of each integral:
\[
 f(z)=\sum_{n=0}^{\infty}(z-z_0)^n
 \underbrace{\frac1{2\pi i}\oint_{C_\rho}
     \frac{f(\xi)}{(\xi-z_0)^{n+1}}\,d\xi}_{=\,f^{(n)}(z_0)/n!\ \text{by \eqref{eq:CIF-deriv}}}.
\]
This is the Taylor series. Its coefficients are the ones found by matching derivatives at the start
of the chapter \eqref{eq:taylor-coef}. The new information is \emph{where} the series equals \(f\):
in every disk around \(z_0\) in which \(f\) is holomorphic.
\par\framebreak\par
\textbf{Consequence: the radius of convergence.} The Taylor series about \(z_0\) converges (and
equals \(f\)) out to the nearest \emph{singularity} of \(f\), a point where \(f\) cannot be made
holomorphic (Section~1.6 classifies them). It cannot converge in a bigger disk: a convergent
power series is holomorphic inside its disk, so it would make \(f\) holomorphic at the
singularity, which is impossible.

\textbf{Consequence: no \(e^{-1/x^2}\) surprise.} For holomorphic functions, ``differentiable'' and
``equal to its Taylor series'' are the same thing. The function \(e^{-1/z^2}\) is not a
counterexample, because it is not even continuous at \(0\): on the imaginary axis \(z=iy\),
\(e^{-1/z^2}=e^{1/y^2}\to\infty\).

\textbf{Why complex numbers matter even for real functions.} The series
\(1/(1+x^2)=1-x^2+x^4-\cdots\) diverges for \(|x|>1\), although \(1/(1+x^2)\) is smooth on the whole
real line. Nothing real explains this. The reason is complex: \(1/(1+z^2)\) has poles at
\(z=\pm i\), at distance \(1\) from \(0\), so the radius of convergence is \(1\).
\end{frame}

\begin{frame}[allowframebreaks=1.00]{Analytic continuation}
\footnotesize

\textbf{Why we want it.} Chapter~2 defines \(\Gamma(z)\) by an integral that converges only for
\(\operatorname{Re}z>0\), and then extends \(\Gamma\) to the left half-plane. Extending a holomorphic
function beyond the region where its first formula works is called \emph{analytic continuation}.
Two questions: how can it be done, and is the result unique?

\textbf{Example: how it can be done.} The geometric series
\[
 f(z)=\frac1{1-z}=1+z+z^2+\cdots
\]
converges only for \(|z|<1\), because of the singularity at \(z=1\). Yet \(1/(1-z)\) is holomorphic at
every other point. The series simply cannot see beyond its disk.

\textbf{Re-expand about a new center.} Pick a point inside the first disk, say \(z=-\tfrac12\), and
expand in powers of \(z+\tfrac12\). Write \(1-z=\tfrac32-(z+\tfrac12)\) and factor out \(\tfrac32\):
\[
    f(z)
    =\frac{1}{\frac32-(z+\frac12)}
    =\frac23\cdot
      \frac{1}{1-\frac23(z+\frac12)}
    =\frac23\sum_{n=0}^{\infty}
      \Bigl[\frac23\Bigl(z+\frac12\Bigr)\Bigr]^n,
\]
using \(1/(1-w)=\sum w^n\) with \(w=\frac23(z+\frac12)\).
\par\framebreak\par
\textbf{Where does the new series converge?} Exactly when \(|w|<1\), that is,
\(|z+\frac12|<\frac32\). This is a new disk \(D_1\) (center \(-\frac12\), radius \(\frac32\)). It overlaps the
first disk \(D_0\) and reaches farther, out to \(z=-2\). On the overlap both series equal
\(1/(1-z)\).

\textbf{Repeat.} Choosing new centers again and again gives a chain of disks
\(D_0\to D_1\to D_2\to\cdots\) (figure on the identity-theorem slide). Each new series agrees with the
previous one on their overlap. In this way \(f\) is carried, step by step, through a larger
region.

\textbf{Why the result is unique.} One might worry that a different chain of disks gives a
different function. The following theorem rules this out.

\begin{block}{Identity theorem}
Let $f,g$ be holomorphic on a connected open set $\Omega$. If they agree on a set with
a limit point inside $\Omega$ (for example, a sequence $z_n\to a$, a segment, or a small
disk), then $f=g$ on all of $\Omega$.
\end{block}
Equivalently (take \(g=0\)): if a holomorphic function on a connected region vanishes on a whole
segment, or on a sequence \(z_n\to a\), it vanishes everywhere. In other words, the zeros of a
holomorphic function that is not identically zero are isolated.
\par\framebreak\par
\textbf{Idea of the proof.} Near any point, a holomorphic function is its Taylor series, so it is
completely fixed by its Taylor coefficients there. Zeros piling up at \(a\) force all those
coefficients to be zero. Then we walk from disk to disk across the region.

\textbf{Proof.} Let \(h=f-g\). It is holomorphic and \(h(z_n)=0\) on a sequence \(z_n\to a\),
\(z_n\ne a\). We show \(h=0\) everywhere.

\textbf{Step 1: \(h\) is zero on a whole disk around \(a\).} Expand \(h\) in its Taylor series at \(a\),
\(h(z)=\sum_k c_k(z-a)^k\). Suppose, to get a contradiction, that not all \(c_k\) are zero, and let
\(c_k\) be the first nonzero one. Then
\[
 h(z)=(z-a)^k\bigl[c_k+c_{k+1}(z-a)+c_{k+2}(z-a)^2+\cdots\bigr].
\]
The bracket tends to \(c_k\neq0\) as \(z\to a\), so it is nonzero for all \(z\) close to \(a\); and
\((z-a)^k\neq0\) for \(z\ne a\). So \(h(z)\ne0\) for every \(z\ne a\) close to \(a\). But the points \(z_n\) come
arbitrarily close to \(a\) and \(h(z_n)=0\): a contradiction. Hence every \(c_k=0\), and \(h=0\) on a disk
\(D_0\) around \(a\).
\par\framebreak\par
\begin{columns}[T]
\column{0.52\textwidth}
\textbf{Step 2: spread the zeros across \(\Omega\).} Take any point \(b\) in \(\Omega\). Since \(\Omega\) is
connected, a path inside \(\Omega\) joins \(a\) to \(b\). Cover the path by finitely many overlapping disks
\(D_0,D_1,\dots,D_N\) inside \(\Omega\) (right).
\begin{itemize}
  \item \(h=0\) on \(D_0\) (Step~1).
  \item \(D_1\) overlaps \(D_0\), so \(h=0\) on an open piece of \(D_1\). That piece contains limit points,
        so Step~1, applied in \(D_1\), gives \(h=0\) on all of \(D_1\).
  \item Continue disk by disk. After \(D_N\), \(h(b)=0\).
\end{itemize}
Since \(b\) was arbitrary, \(h=0\) on \(\Omega\), that is, \(f=g\).
\column{0.44\textwidth}
\begin{center}
    \safeincludegraphics[width=\linewidth]
    {figs/identity-theorem-disks.pdf}
\end{center}
\end{columns}
\par\framebreak\par
\textbf{How we will use it.} Two patterns recur in later chapters.
\begin{itemize}
  \item \textbf{Uniqueness of continuation.} Two continuations of the same function agree on the
        original region, which contains limit points. So they agree wherever both are defined.
        For example, however \(\Gamma\) is continued, the result is the same function.
  \item \textbf{Proving identities.} To prove that two holomorphic functions are equal on a
        region, it is enough to prove it on a segment. Chapter~2 proves
        \(\Gamma(z)\Gamma(1-z)=\pi/\sin(\pi z)\) for real \(z\) in \((0,1)\) with an integral; the
        identity theorem then extends it to every \(z\notin\mathbb Z\), with no further computation.
\end{itemize}
\end{frame}

\begin{frame}[allowframebreaks=1.00]{Laurent series}
\footnotesize

\textbf{Why.} A Taylor series uses only the powers \((z-z_0)^0,(z-z_0)^1,(z-z_0)^2,\dots\), and it
needs \(f\) to be holomorphic at \(z_0\). Near a singularity we also need negative powers
\(1/(z-z_0),\,1/(z-z_0)^2,\dots\). A series with both kinds of powers is a \emph{Laurent series}.

\textbf{Goal.} Show that a function holomorphic on a ring around \(z_0\) always has such an
expansion, and find a formula for its coefficients.

\textbf{Setup.} Let \(f\) be holomorphic in the ring (annulus) \(r<|z-z_0|<R\), and fix \(z\) in it.
Draw two counterclockwise circles around \(z_0\): \(C_1\) of radius \(\rho_1\) outside \(z\), and \(C_2\) of
radius \(\rho_2\) inside \(z\):
\[
 r<\rho_2<|z-z_0|<\rho_1<R .
\]
\textbf{Step 1: Cauchy's formula for the ring between \(C_2\) and \(C_1\).} The boundary of this ring
is \(C_1\) counterclockwise plus \(C_2\) \emph{clockwise} (the region must stay on the left). Joining
the two circles by a slit, as in the proof of \eqref{eq:CIF}, gives
\[
 f(z)=\frac1{2\pi i}\oint_{C_1}\frac{f(\xi)\,d\xi}{\xi-z}-\frac1{2\pi i}\oint_{C_2}\frac{f(\xi)\,d\xi}{\xi-z},
\]
where the minus sign accounts for traversing \(C_2\) clockwise.
\par\framebreak\par
\textbf{Step 2: on \(C_1\), expand as for Taylor.} There \(|z-z_0|<|\xi-z_0|\), so, exactly as on
the Taylor slide,
\[
 \frac1{\xi-z}=\sum_{n=0}^{\infty}\frac{(z-z_0)^n}{(\xi-z_0)^{n+1}}
 \qquad(\text{nonnegative powers of } z-z_0).
\]
\textbf{Step 3: on \(C_2\), expand the other way.} Now \(|\xi-z_0|<|z-z_0|\), so we must factor out
\(z-z_0\) instead. Write \(\xi-z=-\bigl[(z-z_0)-(\xi-z_0)\bigr]\):
\[
 \frac1{\xi-z}=\frac{-1}{(z-z_0)\bigl[1-\frac{\xi-z_0}{z-z_0}\bigr]}
 =-\frac{1}{z-z_0}\sum_{n=0}^{\infty}\Bigl(\frac{\xi-z_0}{z-z_0}\Bigr)^n
 =-\sum_{n=0}^{\infty}\frac{(\xi-z_0)^n}{(z-z_0)^{n+1}} .
\]
These are the negative powers \((z-z_0)^{-1},(z-z_0)^{-2},\dots\). (The \(-\) here cancels the \(-\) in
Step~1.)

\textbf{Step 4: result.} Integrating term by term, \(f\) has a \emph{Laurent expansion},
convergent in the whole ring:
\begin{equation}
 f(z)=\sum_{n=-\infty}^{\infty}a_n(z-z_0)^n
 =\underbrace{\sum_{n=0}^{\infty}a_n(z-z_0)^n}_{\text{regular part}}
  +\underbrace{\sum_{n=1}^{\infty}a_{-n}(z-z_0)^{-n}}_{\text{principal part}}.
 \label{eq:laurent}
\end{equation}
\par\framebreak\par
\textbf{Step 5: a formula for the coefficients.} To isolate one coefficient \(a_k\), divide
\eqref{eq:laurent} by \((z-z_0)^{k+1}\) and integrate counterclockwise around a circle
\(\gamma:|z-z_0|=\rho\) in the ring. Integrating term by term,
\[
 \oint_\gamma\frac{f(z)}{(z-z_0)^{k+1}}\,dz
 =\sum_{n=-\infty}^{\infty}a_n\oint_\gamma(z-z_0)^{n-k-1}\,dz .
\]
By the circle integral \eqref{eq:circ-power}, \(\oint_\gamma(z-z_0)^{n-k-1}dz\) is \(2\pi i\) when the
exponent is \(n-k-1=-1\), that is, \(n=k\), and \(0\) otherwise. So only the term \(n=k\) survives, and
the right side is \(2\pi i\,a_k\). Dividing by \(2\pi i\), for every integer \(k\),
\begin{equation}
 a_k=\frac{1}{2\pi i}\oint_\gamma
      \frac{f(z)}{(z-z_0)^{k+1}}\,dz.
 \label{eq:laurent-coef}
\end{equation}
If the principal part is empty (\(a_k=0\) for all \(k<0\)), the Laurent series is just a Taylor series.
\end{frame}

% ======================================================================
\section{Singularities}
% ======================================================================

\begin{frame}[allowframebreaks=1.00]{Classifying singularities}
\footnotesize

\textbf{Why.} The residue formulas of Section~1.7 depend on how badly \(f\) misbehaves at a point.
We sort the bad points by their Laurent series.

\textbf{Definitions.} A point \(z_0\) is a \emph{singular point} of \(f\) if \(f\) is not holomorphic
there. It is an \emph{isolated singularity} if \(f\) is holomorphic on a punctured disk
\(0<|z-z_0|<R\) (a disk with its center removed). A punctured disk is the ring of \eqref{eq:laurent}
with \(r=0\), so near an isolated singularity
\[
 f(z)=\sum_{n=-\infty}^{\infty}a_n(z-z_0)^n .
\]
The negative powers decide the type.
\par\framebreak\par
\textbf{Type 1: removable} (no negative powers: \(a_n=0\) for all \(n<0\)). Then \(f\) agrees with a power
series, and setting \(f(z_0)=a_0\) makes \(f\) holomorphic at \(z_0\): the singularity was only apparent.

\emph{Example:} \(\sin z/z\) at \(0\). Divide the sine series \eqref{eq:sin-series} by \(z\):
\(\frac{\sin z}{z}=1-\frac{z^2}{3!}+\frac{z^4}{5!}-\cdots\). No negative powers; set the value at \(0\) to \(1\).

\textbf{Type 2: pole of order \(m\)} (finitely many negative powers, the lowest being \((z-z_0)^{-m}\),
with \(a_{-m}\ne0\)). Equivalently,
\[
 f(z)=\frac{h(z)}{(z-z_0)^m},\qquad h\text{ holomorphic near }z_0,\quad h(z_0)\ne0 .
\]
(Multiplying the Laurent series by \((z-z_0)^m\) removes all negative powers and leaves
\(h(z)=a_{-m}+a_{-m+1}(z-z_0)+\cdots\).) A pole of order \(1\) is called \emph{simple}.

\textbf{Type 3: essential} (infinitely many negative powers). \emph{Example:} put \(1/z\) in place of
\(x\) in the sine series (which converges for every complex argument):
\[
 \sin\Bigl(\frac1z\Bigr)=\frac1z-\frac1{3!\,z^3}+\frac1{5!\,z^5}-\cdots .
\]
\par\framebreak\par
\textbf{Meromorphic functions.} A function that is holomorphic except for isolated poles is called
\emph{meromorphic}.

\textbf{Example: \(1/\sin z\) has simple poles exactly at \(z=n\pi\).} We must find all complex zeros of
\(\sin z\).
\begin{itemize}
  \item The derivation of \eqref{eq:cos-sin-exp} works for complex \(z\) too, so
        \(\sin z=(e^{iz}-e^{-iz})/(2i)\). This is zero exactly when \(e^{iz}=e^{-iz}\), that is, when
        \(e^{2iz}=1\).
  \item Write \(z=x+iy\). Then \(2iz=-2y+2ix\), and by \eqref{eq:exp-uv}
        \(e^{2iz}=e^{-2y}(\cos2x+i\sin2x)\).
  \item For this to equal \(1\): its modulus \(e^{-2y}\) must be \(1\), so \(y=0\); then \(\cos2x=1\) and
        \(\sin2x=0\), so \(2x=2\pi n\), that is, \(x=n\pi\).
\end{itemize}
So the zeros of \(\sin z\) are the real numbers \(n\pi\) only. Each is simple: near \(n\pi\),
\(\sin z\approx\cos(n\pi)\,(z-n\pi)\) (first Taylor term), and \(\cos n\pi=\pm1\ne0\).
\par\framebreak\par
\textbf{Branch points are not isolated singularities.} Going once around \(0\) adds \(2\pi i\) to
\(\log z\). So no single-valued branch of \(\log z\) is holomorphic on a whole punctured disk around
\(0\), and there is no Laurent series there.

\textbf{Behaviour at infinity.} Whether a large arc integral vanishes (Section~1.8) depends on how
\(f\) behaves for large \(|z|\). Classify it in the same way by looking at \(F(w)=f(1/w)\) near \(w=0\).
For example, a polynomial of degree \(m\) has a pole of order \(m\) at infinity.
\end{frame}

% ======================================================================
\section{Residues}
% ======================================================================

\begin{frame}[allowframebreaks=1.00]{The residue theorem}
\footnotesize

\textbf{Why we want it.} Many real integrals in physics can be turned into a contour integral
around a closed loop (Section~1.8). The residue theorem evaluates such a loop integral by looking
at one number at each singularity inside.

\textbf{The residue.} Let \(z_0\) be an isolated singularity, with Laurent series
\(f(z)=\sum_n a_n(z-z_0)^n\) near it. The coefficient \(a_{-1}\) of \(1/(z-z_0)\) is called the
\emph{residue}:
\[
 \operatorname{Res}(f,z_0)=a_{-1}.
\]
\textbf{Why this one coefficient?} Integrate the Laurent series around a small counterclockwise
circle \(C\) about \(z_0\), term by term. By the circle integral \eqref{eq:circ-power}, every power
integrates to \(0\) except \((z-z_0)^{-1}\), which gives \(2\pi i\). So only \(a_{-1}\) survives:
\begin{equation}
 \oint_C f(z)\,dz
 =\sum_{n=-\infty}^{\infty}a_n\oint_C(z-z_0)^n\,dz
 =a_{-1}\cdot2\pi i
 =2\pi i\,\operatorname{Res}(f,z_0).
 \label{eq:res-single}
\end{equation}
The residue is what is ``left over'' after integrating around the point.
\par\framebreak\par
\begin{block}{Residue theorem}
Let $f$ be holomorphic inside and on a positively oriented simple closed contour $\gamma$,
except for finitely many isolated singularities $z_1,\ldots,z_N$ in the interior.  Then
\begin{equation}
 \oint_\gamma f(z)\,dz
 =2\pi i\sum_{k=1}^{N}\operatorname{Res}(f,z_k).
 \label{eq:res-thm}
\end{equation}
\end{block}
\textbf{Why.} Away from the singularities \(f\) is holomorphic, so (by Cauchy's theorem, with slits
as in the proof of \eqref{eq:CIF}) we may shrink \(\gamma\) without changing the integral, until it
becomes a set of small circles, one around each \(z_k\). Each small circle gives
\(2\pi i\operatorname{Res}(f,z_k)\) by \eqref{eq:res-single}. Add them up.

\textbf{So the job reduces to computing residues.} We need practical formulas that avoid writing
out the whole Laurent series.
\par\framebreak\par
\textbf{Formula 1: simple pole.} Near a simple pole,
\(f(z)=\frac{a_{-1}}{z-z_0}+a_0+a_1(z-z_0)+\cdots\). Multiply by \(z-z_0\):
\[
 (z-z_0)f(z)=a_{-1}+a_0(z-z_0)+a_1(z-z_0)^2+\cdots\ \xrightarrow{\ z\to z_0\ }\ a_{-1}.
\]
So \(\operatorname{Res}(f,z_0)=\lim_{z\to z_0}(z-z_0)f(z)\).

\textbf{Formula 2: pole of order \(m\).} Now \(f=h/(z-z_0)^m\) with \(h(z)=(z-z_0)^mf(z)\) holomorphic.
Multiplying the Laurent series by \((z-z_0)^m\),
\[
 h(z)=a_{-m}+a_{-m+1}(z-z_0)+\cdots+a_{-1}(z-z_0)^{m-1}+\cdots .
\]
So \(a_{-1}\) is the coefficient of \((z-z_0)^{m-1}\) in the Taylor series of \(h\), which is
\(h^{(m-1)}(z_0)/(m-1)!\) by \eqref{eq:taylor-coef}:
\begin{equation}
 \operatorname{Res}(f,z_0)
 =\frac{1}{(m-1)!}
  \left.\frac{d^{m-1}}{dz^{m-1}}
  \left[(z-z_0)^m f(z)\right]\right|_{z=z_0}.
 \label{eq:res-order-m}
\end{equation}
For \(m=1\) this is Formula~1; for \(m=2\) it is the first derivative of \((z-z_0)^2f(z)\) at \(z_0\).
\par\framebreak\par
\textbf{Formula 3: a quotient \(p/q\) where \(q\) has a simple zero.} Let \(p,q\) be holomorphic near
\(z_0\), with \(q(z_0)=0\) and \(q'(z_0)\ne0\). The Taylor series of \(q\) starts
\(q(z)=q'(z_0)(z-z_0)+\frac12q''(z_0)(z-z_0)^2+\cdots\), so \(q(z)/(z-z_0)\to q'(z_0)\). By Formula~1,
\begin{equation}
 \operatorname{Res}\left(\frac pq,z_0\right)
 =\lim_{z\to z_0}\frac{(z-z_0)\,p(z)}{q(z)}
 =\lim_{z\to z_0}\frac{p(z)}{q(z)/(z-z_0)}
 =\frac{p(z_0)}{q'(z_0)}.
 \label{eq:res-ratio}
\end{equation}
(If also \(p(z_0)=0\), the singularity is removable and the formula correctly gives \(0\).)
\par\framebreak\par
\textbf{Example 1: \(1/(z^2+1)\) at \(z=i\).} Use Formula~3 with \(p=1\), \(q=z^2+1\), \(q'(z)=2z\):
\[
 \operatorname{Res}\Bigl(\frac1{z^2+1},i\Bigr)=\frac{1}{q'(i)}=\frac1{2i}.
\]
\textbf{Example 2: \(1/(z^2+1)^2\) at \(z=i\).} Factor \(z^2+1=(z-i)(z+i)\), so
\(\frac1{(z^2+1)^2}=\frac1{(z-i)^2(z+i)^2}\): a pole of order \(2\) at \(i\). Use Formula~2 with \(m=2\).
First, \((z-i)^2f(z)=\frac1{(z+i)^2}\). Then differentiate once and set \(z=i\):
\[
 \operatorname{Res}
 =\left.\frac{d}{dz}\frac1{(z+i)^2}\right|_{z=i}
  =\left.\frac{-2}{(z+i)^3}\right|_{z=i}
 =\frac{-2}{(2i)^3}
  =\frac{-2}{-8i}
  =\frac1{4i},
\]
using \((2i)^3=8i^3=-8i\).
\par\framebreak\par
\textbf{Two quick contour integrals on \(|z|=1\).}
\begin{itemize}
  \item \(\oint_{|z|=1}\frac{\sin z}{z^2}\,dz\). Divide the sine series by \(z^2\):
        \(\frac{\sin z}{z^2}=\frac{z-z^3/3!+\cdots}{z^2}=\frac1z-\frac{z}{3!}+\cdots\). The coefficient of
        \(1/z\) is \(1\), so the residue at \(0\) is \(1\), and by \eqref{eq:res-single} the integral is
        \(2\pi i\).
  \item \(\oint_{|z|=1}\frac{z^2}{\sin z}\,dz\). Near \(0\),
        \(\frac{z^2}{\sin z}=\frac{z^2}{z(1-z^2/3!+\cdots)}=\frac{z}{1-z^2/3!+\cdots}\), which tends to
        \(0\): no negative powers, so the singularity at \(0\) is removable. The other zeros \(n\pi\) of
        \(\sin z\) lie outside the unit circle. So the integrand is holomorphic inside, and by Cauchy's
        theorem the integral is \(0\).
\end{itemize}
\end{frame}



% ======================================================================
\section{Residue calculus}
% ======================================================================

\begin{frame}[allowframebreaks=1.00]{Real integrals by a semicircular contour}
\footnotesize

\textbf{Why go complex.} An antiderivative of \(1/(1+x^4)\) is messy, and \(\cos x/(1+x^2)\) has no
elementary antiderivative at all. Yet \(\int_{-\infty}^\infty\) of both is easy with residues: close
the real line into a loop, and the residue theorem does the work.

\textbf{The recipe.}
\begin{enumerate}
  \item \emph{Close the path.} Go along the real axis from \(-R\) to \(R\), then back along the upper
        semicircle \(C_R\): \(z=Re^{i\theta}\), \(0\le\theta\le\pi\) (figure below). Call the closed path
        \(\Gamma_R\). It splits into two pieces:
        \begin{equation}
         \oint_{\Gamma_R}f(z)\,dz
         =\int_{-R}^{R}f(x)\,dx+\int_{C_R}f(z)\,dz.
         \label{eq:semicircle-decomposition}
        \end{equation}
  \item \emph{Left side:} by the residue theorem, \(2\pi i\times\)(sum of the residues inside).
  \item \emph{Show the arc integral tends to \(0\)} as \(R\to\infty\) (method below).
  \item \emph{Let \(R\to\infty\).} The real-axis piece becomes the integral we want.
\end{enumerate}
\par\framebreak\par
\textbf{How to show the arc vanishes: the \(ML\) estimate.} An integral is at most (largest size of
the integrand) \(\times\) (length of the path). On \(C_R\), let \(M_R\) be the largest value of \(|f|\); the
length is \(\pi R\). So
\[
 \left|\int_{C_R}f(z)\,dz\right|\le\pi R\, M_R .
\]
If \(R\,M_R\to0\), the arc integral vanishes. This happens, for example, when \(|f|\) falls off like
\(1/R^2\).

\textbf{Result of the recipe.} If the arc vanishes, letting \(R\to\infty\) in
\eqref{eq:semicircle-decomposition} gives
\begin{equation}
 \int_{-\infty}^{\infty}f(x)\,dx
 =2\pi i\sum_{\operatorname{Im}z_k>0}\operatorname{Res}(f,z_k),
 \label{eq:semicircle-result}
\end{equation}
the sum running over the poles \(z_k\) of \(f\) in the upper half-plane.

\textbf{Why upward for \(e^{iaz}\), \(a>0\)?} With \(z=x+iy\),
\(|e^{iaz}|=|e^{iax}|\,|e^{-ay}|=e^{-ay}\). Above the axis (\(y>0\)) this is at most \(1\); below the axis it
grows like \(e^{a|y|}\). So close upward, where \(e^{iaz}\) stays small.

\begin{center}
  \safeincludegraphics[width=0.5\linewidth]{figs/semicircle-contour.pdf}
\end{center}
\par\framebreak\par
\textbf{Example 1: \(\int_{-\infty}^{\infty}\frac{\cos x}{1+x^2}\,dx\).} Since \(\cos x=\operatorname{Re}e^{ix}\),
compute \(J=\int_{-\infty}^{\infty}\frac{e^{ix}}{1+x^2}\,dx\) and take its real part at the end.

\emph{Step 1: the arc vanishes.} On \(C_R\) (with \(R>1\)):
\begin{itemize}
  \item numerator: \(|e^{iz}|=e^{-y}\le1\), since \(y\ge0\);
  \item denominator: \(|1+z^2|\ge|z|^2-1=R^2-1\). (Adding \(1\) to a number of size \(R^2\) cannot make
        it smaller than \(R^2-1\); this is the reverse triangle inequality \(|a+b|\ge|a|-|b|\).)
\end{itemize}
So \(M_R\le1/(R^2-1)\), and
\[
 \left|\int_{C_R}\frac{e^{iz}}{1+z^2}\,dz\right|
 \le\frac{\pi R}{R^2-1}\longrightarrow0 .
\]
\emph{Step 2: the poles.} \(1+z^2=0\) at \(z=\pm i\). Only \(z=i\) is in the upper half-plane.
\par\framebreak\par
\emph{Step 3: the residue at \(i\).} Use \eqref{eq:res-ratio} with \(p=e^{iz}\), \(q=1+z^2\), \(q'=2z\):
\[
 \operatorname{Res}\left(\frac{e^{iz}}{1+z^2},i\right)=\frac{p(i)}{q'(i)}=\frac{e^{i\cdot i}}{2i}=\frac{e^{-1}}{2i}.
\]
\emph{Step 4: the answer.} By \eqref{eq:semicircle-result},
\[
 J=2\pi i\cdot\frac{e^{-1}}{2i}=\frac{\pi}{e}.
\]
This is real, so its real part is itself:
\begin{equation}
 \int_{-\infty}^{\infty}\frac{\cos x}{1+x^2}\,dx=\frac{\pi}{e}.
 \label{eq:cos-rational-integral}
\end{equation}
\emph{Check:} the imaginary part, \(\int\sin x/(1+x^2)\,dx\), must be \(0\), and it is, because
\(\sin x/(1+x^2)\) is odd.
\par\framebreak\par
\textbf{Example 2: \(\int_{-\infty}^{\infty}\frac{dx}{1+x^4}\).}

\emph{Step 1: the arc vanishes.} As before, \(|1+z^4|\ge R^4-1\) on \(C_R\), so
\(\bigl|\int_{C_R}\bigr|\le\pi R/(R^4-1)\to0\).

\emph{Step 2: the poles.} Solve \(z^4=-1=e^{i\pi}\). Taking fourth roots, and remembering that
\(e^{i\pi}=e^{i(\pi+2\pi k)}\) for every integer \(k\),
\[
 z_k=e^{i(\pi+2k\pi)/4}:\qquad
 e^{i\pi/4},\quad e^{3i\pi/4},\quad e^{5i\pi/4},\quad e^{7i\pi/4}\qquad(k=0,1,2,3).
\]
In the upper half-plane (angle between \(0\) and \(\pi\)) are \(z_0=e^{i\pi/4}\) and \(z_1=e^{3i\pi/4}\).

\emph{Step 3: the residues.} Use \eqref{eq:res-ratio} with \(p=1\), \(q=1+z^4\), \(q'=4z^3\). Multiply top and
bottom by \(z_j\), then use \(z_j^4=-1\):
\[
 \operatorname{Res}\Bigl(\frac1{1+z^4},z_j\Bigr)
  =\frac1{4z_j^3}=\frac{z_j}{4z_j^4}=-\frac{z_j}{4}.
\]
\par\framebreak\par
\emph{Step 4: add them.} By Euler's formula \eqref{eq:euler},
\[
 z_0=e^{i\pi/4}=\frac{1+i}{\sqrt2},\qquad z_1=e^{3i\pi/4}=\frac{-1+i}{\sqrt2},\qquad
 z_0+z_1=\frac{2i}{\sqrt2}=i\sqrt2 .
\]
So the sum of the residues is \(-\frac{z_0+z_1}{4}=-\frac{i\sqrt2}{4}\).

\emph{Step 5: the answer.} By \eqref{eq:semicircle-result}, and \(i\cdot(-i)=1\),
\begin{equation}
 \int_{-\infty}^{\infty}\frac{dx}{1+x^4}
  =2\pi i\Bigl(-\frac{i\sqrt2}{4}\Bigr)
  =\frac{2\pi\sqrt2}{4}
  =\frac{\pi}{\sqrt2}.
 \label{eq:one-plus-x4}
\end{equation}
\end{frame}

\begin{frame}[allowframebreaks=1.00]{Keyhole contour: $\int_0^\infty x^a/(1+x)^2\,dx$}
\footnotesize

\textbf{Goal.} Evaluate
\[
 I(a)=\int_0^\infty\frac{x^a}{(1+x)^2}\,dx,
 \qquad -1<a<1.
\]
\textbf{Why \(-1<a<1\)?} Near \(x=0\) the integrand behaves like \(x^a\), integrable when \(a>-1\). For large
\(x\) it behaves like \(x^a/x^2=x^{a-2}\), integrable when \(a-2<-1\), that is, \(a<1\).

\textbf{The difficulty.} The integral runs only over \((0,\infty)\), not the whole real line, so the
semicircle recipe does not apply. Also, \(z^a\) has a branch point at \(0\).

\textbf{The idea: use the branch cut.} Put the branch cut of \(z^a\) along the positive real axis,
exactly where we integrate. Then run a path out along the top side of the cut and back along the
bottom side. For an ordinary function the two sides would cancel. But \(z^a\) is different on the
two sides (by a factor \(e^{2\pi ia}\)), so they do not cancel. What is left over is a multiple of
\(I(a)\), and the residue theorem computes it.
\par\framebreak\par
\textbf{Step 1: choose the branch.} Define
\[
 z^a=e^{a\ell(z)},\qquad
 \ell(z)=\log|z|+i\arg z,\qquad 0<\arg z<2\pi .
\]
This is like the principal logarithm, but the angle runs from \(0\) to \(2\pi\), so the cut is the
positive real axis.

\textbf{Step 2: the keyhole path.} Counterclockwise: out along the top of the cut from \(\varepsilon\) to
\(R\), around the big circle \(C_R\), back along the bottom of the cut from \(R\) to \(\varepsilon\), and around
the small circle \(C_\varepsilon\) (clockwise). The small circle keeps the path away from the branch point
\(0\), where there is no residue to use.

\begin{center}
  \safeincludegraphics[height=0.42\textheight]{figs/keyhole-contour.pdf}
\end{center}
\par\framebreak\par
\textbf{Step 3: the top side.} There \(\arg z\to0\), so \(z=x\) and \(z^a=x^a\). The contribution is
\[
 I_{\varepsilon,R}(a)=\int_\varepsilon^R\frac{x^a}{(1+x)^2}\,dx .
\]
\textbf{Step 4: the bottom side.} There \(\arg z\to2\pi\) (we went once around), so \(z=xe^{2\pi i}\) and
\(\ell(z)=\log x+2\pi i\). Hence
\[
 z^a=e^{a(\log x+2\pi i)}=x^a\,e^{2\pi ia}.
\]
The direction is from \(R\) back to \(\varepsilon\), which gives a minus sign:
\[
 \int_R^\varepsilon\frac{x^a\,e^{2\pi ia}}{(1+x)^2}\,dx
 =-e^{2\pi ia}\,I_{\varepsilon,R}(a).
\]
\textbf{The two sides together} give \(\bigl(1-e^{2\pi ia}\bigr)\,I_{\varepsilon,R}(a)\). Because \(a\) is not an
integer, \(e^{2\pi ia}\neq1\), so this is not zero.
\par\framebreak\par
\textbf{Step 5: the two circles vanish} (by the \(ML\) estimate).
\begin{itemize}
  \item \emph{Big circle} \(|z|=R>1\): \(|z^a|=R^a\); \(|1+z|\ge R-1\) (reverse triangle inequality); length
        \(2\pi R\). So
        \[
         \left|\int_{C_R}\frac{z^a}{(1+z)^2}\,dz\right|
         \le \frac{2\pi R\cdot R^{a}}{(R-1)^2}\approx 2\pi R^{a-1}\longrightarrow0
         \qquad(\text{since }a<1).
        \]
  \item \emph{Small circle} \(|z|=\varepsilon<1\): \(|z^a|=\varepsilon^a\); \(|1+z|\ge1-\varepsilon\); length
        \(2\pi\varepsilon\). So
        \[
         \left|\int_{C_\varepsilon}\frac{z^a}{(1+z)^2}\,dz\right|
         \le \frac{2\pi\varepsilon\cdot\varepsilon^{a}}{(1-\varepsilon)^2}\longrightarrow0
         \qquad(\text{since }a>-1).
        \]
\end{itemize}
\par\framebreak\par
\textbf{Step 6: the residue.} The only singularity inside is \(z=-1\), a double pole. Use
\eqref{eq:res-order-m} with \(m=2\): first \((z+1)^2\cdot\frac{z^a}{(1+z)^2}=z^a\); then differentiate,
\((z^a)'=a\,z^{a-1}\) (with \(z^{a-1}=e^{(a-1)\ell(z)}\)); then set \(z=-1\). On our branch \(\arg(-1)=\pi\),
so \(\ell(-1)=i\pi\) and
\[
 (-1)^{a-1}=e^{(a-1)i\pi}=e^{i\pi a}\,e^{-i\pi}=-e^{i\pi a}.
\]
Therefore
\[
 \operatorname{Res}\left(\frac{z^a}{(1+z)^2},-1\right)=a\cdot\bigl(-e^{i\pi a}\bigr)=-a\,e^{i\pi a}.
\]
\textbf{Step 7: the residue theorem.} Let \(R\to\infty\) and \(\varepsilon\to0\). The circles drop out, and
\(I_{\varepsilon,R}\to I(a)\):
\[
 \bigl(1-e^{2\pi ia}\bigr)\,I(a)=2\pi i\cdot\bigl(-a\,e^{i\pi a}\bigr).
\]
\par\framebreak\par
\textbf{Step 8: simplify.} Factor \(e^{i\pi a}\) out of \(1-e^{2\pi ia}\), then use
\(e^{i\theta}-e^{-i\theta}=2i\sin\theta\) from \eqref{eq:cos-sin-exp}:
\[
 1-e^{2\pi ia}=e^{i\pi a}\bigl(e^{-i\pi a}-e^{i\pi a}\bigr)=e^{i\pi a}\cdot\bigl(-2i\sin(\pi a)\bigr).
\]
Divide Step~7 by this. The factors \(-2i\,e^{i\pi a}\) cancel:
\[
 I(a)=\frac{-2\pi i\,a\,e^{i\pi a}}{-2i\,e^{i\pi a}\sin(\pi a)}=\frac{\pi a}{\sin(\pi a)}.
\]
\begin{equation}
 \int_0^\infty\frac{x^a}{(1+x)^2}\,dx
 =\frac{\pi a}{\sin(\pi a)},\qquad -1<a<1.
 \label{eq:keyhole-xa}
\end{equation}
\textbf{Check} at \(a=0\): the right side tends to \(1\) (since \(\sin(\pi a)\approx\pi a\)), and indeed
\(\int_0^\infty dx/(1+x)^2=\bigl[-1/(1+x)\bigr]_0^\infty=1\).
\par\framebreak\par
\textbf{Companion integral, used in Chapter~2.} The reflection formula
\(\Gamma(a)\Gamma(1-a)=\pi/\sin(\pi a)\) rests on
\begin{equation}
 \int_0^\infty\frac{x^{-a}}{1+x}\,dx
 =\frac{\pi}{\sin(\pi a)},\qquad 0<a<1.
 \label{eq:keyhole2}
\end{equation}
Same method with \(z^{-a}/(1+z)\):
\begin{itemize}
  \item \emph{Two sides of the cut:} the bottom side now carries \(e^{-2\pi ia}\), so they give
        \((1-e^{-2\pi ia})\int_0^\infty\frac{x^{-a}}{1+x}dx\).
  \item \emph{Circles:} bounded by \(2\pi R^{1-a}/(R-1)\) and \(2\pi\varepsilon^{1-a}/(1-\varepsilon)\); both tend to
        \(0\) because \(0<a<1\).
  \item \emph{Residue} at the simple pole \(-1\) (Formula~1): the numerator there,
        \((-1)^{-a}=e^{-a\ell(-1)}=e^{-i\pi a}\).
\end{itemize}
So \((1-e^{-2\pi ia})\int_0^\infty\frac{x^{-a}}{1+x}dx=2\pi i\,e^{-i\pi a}\). Factor
\(1-e^{-2\pi ia}=e^{-i\pi a}(e^{i\pi a}-e^{-i\pi a})=2i\,e^{-i\pi a}\sin(\pi a)\) and divide: the result is
\(2\pi i\,e^{-i\pi a}/\bigl(2i\,e^{-i\pi a}\sin(\pi a)\bigr)=\pi/\sin(\pi a)\).
\end{frame}

\begin{frame}[allowframebreaks=1.00]{Jordan's lemma and $\int \sin x/x\,dx$}
\footnotesize

\textbf{Goal.} Evaluate \(\int_{-\infty}^{\infty}\frac{\sin x}{x}\,dx\).

\textbf{Plan.} Since \(\sin x=\operatorname{Im}e^{ix}\), integrate \(e^{iz}/z\) around an upper semicircle,
as in the semicircle recipe. Two new obstacles appear:
\begin{itemize}
 \item the pole \(z=0\) lies \emph{on} the real axis, right on our path. We go around it on a
       small semicircle (Tool~1);
 \item on the big arc, \(|1/z|=1/R\), so the \(ML\) bound is \(\pi R\cdot\frac1R=\pi\), which does not
       tend to \(0\). We need a sharper estimate (Tool~2, Jordan's lemma).
\end{itemize}
\par\framebreak\par
\textbf{Tool 1: a small detour around a pole on the path.} Let \(z_0\) be a simple pole on the real
axis. Near it, \(f(z)=\frac{c}{z-z_0}+g(z)\), where \(c=\operatorname{Res}(f,z_0)\) and \(g\) is holomorphic,
hence bounded, near \(z_0\). Go around \(z_0\) on the upper half-circle \(C_\varepsilon\) of radius \(\varepsilon\),
from left to right (that is, clockwise): \(z=z_0+\varepsilon e^{i\theta}\) with \(\theta\) going from \(\pi\) down
to \(0\).
\begin{itemize}
  \item The pole term: \(z-z_0=\varepsilon e^{i\theta}\) and \(dz=i\varepsilon e^{i\theta}d\theta\), so
        \(\frac{c\,dz}{z-z_0}=ic\,d\theta\), and
        \[
         \int_{C_\varepsilon}\frac{c}{z-z_0}\,dz=ic\int_\pi^0d\theta=ic\,(0-\pi)=-i\pi c .
        \]
  \item The bounded part \(g\): at most (its bound)\(\times\)(length \(\pi\varepsilon\)), which tends to \(0\).
\end{itemize}
So, as \(\varepsilon\to0\), the detour contributes \(-i\pi\) times the residue (half of the full
\(-2\pi i\), since it is half a circle, traversed clockwise):
\begin{equation}
 \int_{C_\varepsilon}f(z)\,dz\longrightarrow
 -i\pi\operatorname{Res}(f,z_0).\label{eq:indent}
\end{equation}
\par\framebreak\par
\textbf{Tool 2: why the big arc still vanishes.} On the arc \(z=Re^{i\theta}\), \(0\le\theta\le\pi\),
\[
 |e^{iaz}|=|e^{iaR\cos\theta}|\cdot|e^{-aR\sin\theta}|=e^{-aR\sin\theta}\qquad(a>0).
\]
This is extremely small except near the two ends \(\theta\approx0\) and \(\theta\approx\pi\), where \(\sin\theta\)
is small. How big are those end pieces? \(e^{-aR\sin\theta}\) is not small only while
\(aR\sin\theta\lesssim1\), that is, \(\theta\lesssim1/(aR)\). On a circle of radius \(R\), that is an arc of
length about \(R\cdot\frac1{aR}=\frac1a\), whatever \(R\) is. So the integral is roughly
\(M_R\times\frac1a\), not \(M_R\times\pi R\), and it vanishes when \(M_R\to0\).

\textbf{Jordan's lemma.} \emph{Let \(a>0\), let \(C_R\) be the upper semicircle of radius \(R\), and let
\(M_R=\max_{z\in C_R}|f(z)|\). If \(M_R\to0\) as \(R\to\infty\), then}
\[
 \int_{C_R}f(z)\,e^{iaz}\,dz\longrightarrow0.
\]
\par\framebreak\par
\textbf{Proof} (making the picture precise).
\begin{enumerate}
  \item Bound by size times length, with \(|dz|=R\,d\theta\):
        \(\bigl|\int_{C_R}fe^{iaz}dz\bigr|\le M_RR\int_0^\pi e^{-aR\sin\theta}d\theta\).
  \item By symmetry, \(\sin(\pi-\theta)=\sin\theta\), so the halves \([0,\frac\pi2]\) and \([\frac\pi2,\pi]\) give the
        same: \(\int_0^\pi=2\int_0^{\pi/2}\).
  \item On \([0,\frac\pi2]\) the graph of \(\sin\theta\) lies above its chord from \((0,0)\) to
        \((\frac\pi2,1)\), so \(\sin\theta\ge2\theta/\pi\), and \(e^{-aR\sin\theta}\le e^{-2aR\theta/\pi}\).
  \item Integrate the exponential exactly.
\end{enumerate}
\begin{align*}
 \left|\int_{C_R}f(z)e^{iaz}\,dz\right|
 &\le2M_R R\int_0^{\pi/2}e^{-2aR\theta/\pi}\,d\theta
 =2M_R R\left[-\frac{\pi}{2aR}e^{-2aR\theta/\pi}\right]_0^{\pi/2}\\
 &=2M_R R\cdot\frac{\pi}{2aR}\bigl(1-e^{-aR}\bigr)
 \le\frac{\pi M_R}{a}\longrightarrow0.
\end{align*}
This matches the picture: \(M_R\) times a length of order \(1/a\).
\par\framebreak\par
\textbf{The Dirichlet integral.} Take the closed path: real axis from \(-R\) to \(-\varepsilon\), small
semicircle \(C_\varepsilon\) over \(0\), real axis from \(\varepsilon\) to \(R\), big semicircle \(C_R\). It encloses no
pole (the pole \(0\) has been avoided), so by Cauchy's theorem
\[
 \int_{-R}^{-\varepsilon}\frac{e^{ix}}x\,dx
 +\int_{C_\varepsilon}\frac{e^{iz}}z\,dz
 +\int_{\varepsilon}^{R}\frac{e^{ix}}x\,dx
 +\int_{C_R}\frac{e^{iz}}z\,dz=0 .
\]
Now let \(R\to\infty\) and \(\varepsilon\to0\), one piece at a time:
\begin{itemize}
 \item \emph{Big arc:} \(f=1/z\) has \(M_R=1/R\to0\), so by Jordan's lemma (\(a=1\)) it tends to \(0\).
 \item \emph{Small arc:} \(\operatorname{Res}(e^{iz}/z,0)=\lim_{z\to0}z\cdot\frac{e^{iz}}z=1\), so by
       \eqref{eq:indent} it tends to \(-i\pi\).
 \item \emph{Real pieces:} their sum, taken symmetrically around \(0\), is called the \emph{Cauchy
       principal value}, written \(\operatorname{PV}\int_{-\infty}^\infty\). (The symmetric limit is
       needed because the real part \(\cos x/x\) blows up like \(1/x\) at \(0\); the two sides cancel.)
\end{itemize}
\par\framebreak\par
\textbf{Result.} The limit gives \(\operatorname{PV}\int_{-\infty}^{\infty}\frac{e^{ix}}x\,dx-i\pi=0\), so
\begin{equation}
 \operatorname{PV}\int_{-\infty}^{\infty}\frac{e^{ix}}x\,dx=i\pi.
 \label{eq:pv-exponential-over-x}
\end{equation}
Write \(e^{ix}=\cos x+i\sin x\) and take imaginary parts. The sine part needs no principal value,
because \(\sin x/x\to1\) at \(x=0\):
\begin{equation}
 \int_{-\infty}^{\infty}\frac{\sin x}{x}\,dx=\pi.
 \label{eq:dirichlet-integral}
\end{equation}
(The real part says \(\operatorname{PV}\int\cos x/x\,dx=0\), as it must, since \(\cos x/x\) is odd.)
\end{frame}

\begin{frame}[allowframebreaks=1.00]{Sampling kernels: writing a sum as an integral}
\footnotesize
\textbf{Why.} Residue calculus also works in reverse. A \emph{sum} \(\sum_n f(n)\) can be written as a
contour integral, and the integral can then be evaluated from the few residues of \(f\) itself.
This gives closed forms for sums such as \(\sum 1/(n^2+a^2)\).

\textbf{The idea.} Find a function \(K(z)\) that has a simple pole with residue \(1\) at every
integer. Then \(f(z)K(z)\) has residue \(f(n)\) at each integer \(n\), and the residue theorem adds up
exactly the values \(f(n)\).

\textbf{The kernel \(K(z)=\pi\cot\pi z=\dfrac{\pi\cos\pi z}{\sin\pi z}\).}
\begin{itemize}
  \item It has a simple pole at each integer \(n\), because \(\sin\pi z=0\) exactly there.
  \item Its residue there is \(1\): use \eqref{eq:res-ratio} with \(p=\pi\cos\pi z\), \(q=\sin\pi z\),
        \(q'=\pi\cos\pi z\). Then \(p(n)/q'(n)=\pi\cos n\pi/(\pi\cos n\pi)=1\).
  \item So \(\operatorname{Res}(f\cdot\pi\cot\pi z,\,n)=f(n)\cdot1=f(n)\) when \(f\) is holomorphic at \(n\).
        (\(K\) ``samples'' \(f\) at the integers.)
\end{itemize}
\textbf{Result.} If a contour \(C\) encloses the integers \(n_1,\dots,n_2\), and \(f\) has no poles inside,
the residue theorem \eqref{eq:res-thm} gives
\[
 \sum_{m=n_1}^{n_2}f(m)
 =\frac{1}{2\pi i}\oint_C f(z)\,\pi\cot\pi z\,dz .
\]
If \(f\) does have poles inside \(C\), the residues of \(f\,\pi\cot\pi z\) at those poles are added too.
\par\framebreak\par
\textbf{Example: \(\displaystyle\sum_{n=-\infty}^{\infty}\frac1{n^2+a^2}\), \(a>0\).} Take \(f(z)=1/(z^2+a^2)\). Let
\(C_N\) be the square with corners \((N+\tfrac12)(\pm1\pm i)\). It passes midway between integers, and it
encloses \(n=-N,\dots,N\) and the poles \(\pm ia\) of \(f\) (for large \(N\)).

\textbf{Step 1: the integral over \(C_N\) tends to \(0\).}
\begin{itemize}
  \item \(\cot\pi z\) stays bounded on \(C_N\), say \(|\cot\pi z|\le2\). (Intuition: the square stays away
        from the poles of \(\cot\) by passing midway between integers; and far from the real axis
        \(\cot\pi z\) tends to the constants \(\mp i\). We accept the bound.)
  \item On \(C_N\), \(|z|\ge N+\tfrac12\), so \(|z^2+a^2|\ge(N+\tfrac12)^2-a^2\).
  \item The length of the square is \(4\times2(N+\tfrac12)=8(N+\tfrac12)\).
\end{itemize}
By the \(ML\) estimate,
\[
 \left|\oint_{C_N}\frac{\pi\cot\pi z}{z^2+a^2}\,dz\right|
 \le8\Bigl(N+\tfrac12\Bigr)\cdot\frac{2\pi}{(N+\tfrac12)^2-a^2}\longrightarrow0 ,
\]
since the bound behaves like \(16\pi/N\).
\par\framebreak\par
\textbf{Step 2: the residues at the poles of \(f\), \(z_0=\pm ia\).} Use \eqref{eq:res-ratio} with
\(p=\pi\cot\pi z\), \(q=z^2+a^2\), \(q'=2z\): the residue is \(\pi\cot(\pi z_0)/(2z_0)\). We need \(\cot\) at an
imaginary argument. From \eqref{eq:cos-sin-exp} with \(\theta=iy\) (valid for complex arguments):
\[
 \cos(iy)=\frac{e^{-y}+e^{y}}2=\cosh y,\qquad
 \sin(iy)=\frac{e^{-y}-e^{y}}{2i}=\frac{-2\sinh y}{2i}=i\sinh y,
\]
so \(\cot(iy)=\frac{\cosh y}{i\sinh y}=-i\coth y\), where \(\coth=\cosh/\sinh\).
\begin{itemize}
  \item At \(z_0=ia\): \(\dfrac{\pi\cdot(-i\coth\pi a)}{2ia}=-\dfrac{\pi\coth\pi a}{2a}\).
  \item At \(z_0=-ia\): \(\cot\) is odd, so \(\cot(-i\pi a)=+i\coth\pi a\), and
        \(\dfrac{\pi\cdot(i\coth\pi a)}{-2ia}=-\dfrac{\pi\coth\pi a}{2a}\).
\end{itemize}
Together: \(-\pi\coth(\pi a)/a\).
\par\framebreak\par
\textbf{Step 3: combine.} The residue theorem on \(C_N\) says
\(\frac1{2\pi i}\oint_{C_N}=\sum_{n=-N}^Nf(n)+\bigl(-\pi\coth(\pi a)/a\bigr)\). As \(N\to\infty\) the left side
tends to \(0\) (Step~1), so
\[
 \sum_{n=-\infty}^{\infty}\frac1{n^2+a^2}=\frac{\pi}{a}\coth(\pi a).
\]
\textbf{Check: let \(a\to0\).} Both sides blow up like \(1/a^2\); compare the next terms.
\begin{itemize}
  \item Left: the \(n=0\) term is \(1/a^2\); the others give
        \(2\sum_{n\ge1}\frac1{n^2+a^2}\to2\sum_{n\ge1}\frac1{n^2}\).
  \item Right: expand \(\coth x=\frac{\cosh x}{\sinh x}=\frac{1+x^2/2+\cdots}{x+x^3/6+\cdots}
        =\frac1x\cdot\frac{1+x^2/2+\cdots}{1+x^2/6+\cdots}\). Using \(\frac1{1+x^2/6}\approx1-\frac{x^2}6\),
        this is \(\frac1x\bigl(1+\frac{x^2}2-\frac{x^2}6\bigr)=\frac1x+\frac x3\). With \(x=\pi a\):
        \(\frac\pi a\coth\pi a=\frac{1}{a^2}+\frac{\pi^2}{3}+\cdots\).
\end{itemize}
Matching: \(2\sum_{n\ge1}1/n^2=\pi^2/3\), so \(\sum_{n\ge1}1/n^2=\pi^2/6\), the Basel value used in
Chapters~2 and~7.
\par\framebreak\par
\textbf{Other sample sets.} The same residue computation with \eqref{eq:res-ratio} gives other kernels:
\begin{itemize}
 \item \(\pi\csc\pi z=\pi/\sin\pi z\): residue \(\pi/(\pi\cos n\pi)=(-1)^n\) at \(z=n\). It samples the
       alternating sum \(\sum(-1)^m f(m)\).
 \item \(\pi/\cos\pi z\): poles at the half-integers \(z=n+\tfrac12\), with residue
       \(\pi/\bigl(-\pi\sin(\pi n+\tfrac\pi2)\bigr)=-1/\cos(\pi n)=(-1)^{n+1}\). It samples
       \(\sum(-1)^{n+1}f(n+\tfrac12)\).
\end{itemize}
\end{frame}

% ======================================================================
\section{Method of steepest descent}
% ======================================================================

\begin{frame}[allowframebreaks=1.00]{The steepest-descent approximation}
\footnotesize

\textbf{Goal.} Many formulas for special functions at large arguments (Stirling, large-\(x\)
Bessel, Airy) come from integrals of the form
\begin{equation}
  I(s)=\int_\Gamma g(z)\,e^{s f(z)}\,dz,
  \qquad s\to+\infty,
  \label{eq:sd-int}
\end{equation}
with \(f,g\) holomorphic near the path \(\Gamma\). We want a simple approximate formula for \(I(s)\) when
the parameter \(s\) is large.

\textbf{The basic idea.} Write \(f=u+iv\) with \(u,v\) real. Then \(|e^{sf}|=e^{su}\). For large \(s\), the
factor \(e^{su}\) is enormously larger where \(u\) is largest than anywhere else. (If \(u\) is bigger by
\(1\) at one point, \(e^{su}\) is bigger there by a factor \(e^s\).) So only a small neighbourhood of the
highest point of \(u\) on the path matters, and near that point we can replace \(f\) by its
Taylor expansion. The result is a Gaussian integral.

\textbf{The complication.} Since \(f\) is holomorphic, we may deform \(\Gamma\) (Cauchy's theorem), as long
as we cross no singularity. Which path should we use? The next steps answer this.
\par\framebreak\par
\textbf{Step 1: the important points are saddle points.} We look for \(z_0\) with \(f'(z_0)=0\) and
\(f''(z_0)\ne0\).
\begin{columns}[T]
\begin{column}{0.62\linewidth}
\begin{itemize}
  \item By Cauchy--Riemann, \(f'=u_x+iv_x=u_x-iu_y\). So \(f'(z_0)=0\) means \(u_x=u_y=0\): \(z_0\) is a flat
        point of the landscape \(u\).
  \item It is never a hilltop. From \eqref{eq:CR}, \(u_x=v_y\) gives \(u_{xx}=v_{yx}\), and \(u_y=-v_x\)
        gives \(u_{yy}=-v_{xy}\). So \(u_{xx}+u_{yy}=0\): if \(u\) curves down in one direction, it
        curves up in the other. The flat point is a mountain pass (a saddle).
\end{itemize}
\textbf{Picture} \(u\) as terrain (right). A path joining two valleys must cross the ridge between
them. The best crossing is through the pass \(z_0\): there the highest point of the route is as low
as possible. From the pass, the path of steepest descent runs down into both valleys.
\end{column}
\begin{column}{0.34\linewidth}
\centering
\safeincludegraphics[width=\linewidth,height=0.42\textheight,keepaspectratio]{figs/saddle-level-curves.pdf}
\end{column}
\end{columns}
\par\framebreak\par
\textbf{Step 2: the direction of steepest descent.} Near \(z_0\), Taylor's theorem gives (no linear
term, because \(f'(z_0)=0\))
\begin{equation}
  f(z)=f(z_0)+\frac12 f''(z_0)(z-z_0)^2+O\bigl((z-z_0)^3\bigr).
  \label{eq:saddle-taylor}
\end{equation}
Write \(f''(z_0)\) and \(z-z_0\) in polar form: \(f''(z_0)=|f''(z_0)|e^{i\theta}\), \(z-z_0=re^{i\phi}\). Then
\[
  \frac12 f''(z_0)(z-z_0)^2
  =\frac12|f''(z_0)|\,r^2\,e^{i(\theta+2\phi)}
  =\frac12|f''(z_0)|\,r^2\bigl[\cos(\theta+2\phi)+i\sin(\theta+2\phi)\bigr].
\]
\begin{itemize}
  \item The real part (which controls the size \(e^{su}\)) drops fastest when
        \(\cos(\theta+2\phi)=-1\).
  \item In that direction the imaginary part \(\sin(\theta+2\phi)\) is \(0\). So the phase of \(e^{sf}\) does
        not rotate there: the integrand does not oscillate, and nothing cancels. This makes the
        approximation clean.
\end{itemize}
So the steepest-descent direction satisfies
\begin{equation}
  \theta+2\phi=\pi\pmod{2\pi},
  \qquad\text{that is}\qquad
  \phi=\frac{\pi-\theta}{2}\pmod{\pi}.
  \label{eq:descent-direction}
\end{equation}
The ``\(\pmod\pi\)'' gives two opposite rays from \(z_0\): the path goes down into both valleys.
(The perpendicular direction, \(\theta+2\phi=0\), is steepest \emph{ascent}.)
\par\framebreak\par
\textbf{Step 3: walk along the descent line.} Deform \(\Gamma\) so that it passes through \(z_0\) along the
descent line. Parametrize the line by a real number \(r\):
\[
  z=z_0+\alpha r,\qquad
  \alpha=e^{i(\pi-\theta)/2}=e^{i\pi/2}e^{-i\theta/2}=i\,e^{-i\theta/2},
  \qquad dz=\alpha\,dr .
\]
Check that this is a descent line: \(\alpha^2=i^2e^{-i\theta}=-e^{-i\theta}\), so
\(f''(z_0)\,\alpha^2=|f''(z_0)|e^{i\theta}\cdot(-e^{-i\theta})=-|f''(z_0)|\), a negative real number. Hence,
along the line,
\[
  s\,f(z)\approx s\,f(z_0)-\frac12\,s\,|f''(z_0)|\,r^2 .
\]
\textbf{Step 4: only a small window matters.} The factor \(\exp\bigl(-\frac12s|f''(z_0)|r^2\bigr)\) is
negligible once \(|r|\) is a few times \(1/\sqrt{s|f''(z_0)|}\). This window shrinks as \(s\) grows. Inside
it, the cubic term \(s\cdot O(r^3)\) is of size \(s\cdot s^{-3/2}=s^{-1/2}\to0\), and \(g(z)\approx g(z_0)\). So
we may replace \(f\) by its quadratic approximation, \(g\) by \(g(z_0)\), and integrate over all \(r\).
\par\framebreak\par
\textbf{Step 5: do the Gaussian integral.} With these replacements
(\(A\sim B\) means \(A/B\to1\)),
\[
  I(s)
  \sim \int_{-\infty}^{\infty} g(z_0)\,
    e^{sf(z_0)}\,e^{-\frac12 s|f''(z_0)|r^2}\,\alpha\,dr
  =g(z_0)\,e^{sf(z_0)}\,\alpha
    \int_{-\infty}^{\infty}e^{-\frac12 s|f''(z_0)|r^2}\,dr .
\]
Use \(\int_{-\infty}^{\infty}e^{-cr^2}\,dr=\sqrt{\pi/c}\) (Chapter~2's Gaussian integral, after
\(r=x/\sqrt c\)) with \(c=\tfrac12 s|f''(z_0)|\):
\(\sqrt{\pi/c}=\sqrt{2\pi/(s|f''(z_0)|)}\). With \(\alpha=i\,e^{-i\theta/2}\):
\begin{equation}
  I(s)\sim
  g(z_0)\,e^{sf(z_0)}\;i\,e^{-i\theta/2}\;
  \sqrt{\frac{2\pi}{s|f''(z_0)|}}.
  \label{eq:sd-result}
\end{equation}
\textbf{Reading the formula.} \(e^{sf(z_0)}\) is the height of the pass; \(\sqrt{2\pi/(s|f''|)}\) is the
width of the peak; \(i\,e^{-i\theta/2}\) is the direction of the path through the pass.
\par\framebreak\par
\textbf{Practical remarks.}
\begin{itemize}
  \item \emph{Orientation.} Going along the line in the opposite direction flips the sign of the
        answer. Choose the direction in which the original path crosses the pass.
  \item \emph{Several saddles} of the same height: add their contributions.
  \item \emph{Endpoint maximum.} If the highest point is an endpoint of the path instead of a
        saddle, the answer is smaller: for example, integrating by parts once,
        \(\int_0^\infty g(t)e^{-st}dt\approx g(0)/s\), which falls off like \(1/s\) rather than
        \(1/\sqrt s\).
\end{itemize}
\end{frame}

\begin{frame}[allowframebreaks=1.00]{Real special cases: Laplace method and stationary phase}
\footnotesize

The same local-Gaussian idea has two real-variable versions that we use later.

\textbf{1. Laplace's method (a decaying exponential).} Let \(\varphi\) be real with a single
maximum at an interior point \(x_0\) (\(\varphi'(x_0)=0\), \(\varphi''(x_0)<0\)), and \(g(x_0)\ne0\). Then
\begin{equation}
  \int_a^b g(x)\,e^{s\varphi(x)}\,dx
  \sim
  g(x_0)\,e^{s\varphi(x_0)}
  \sqrt{\frac{2\pi}{s\bigl|\varphi''(x_0)\bigr|}},
  \qquad s\to+\infty.
  \label{eq:laplace-real}
\end{equation}
\textbf{Why.} Near the maximum, \(\varphi(x)\approx\varphi(x_0)-\frac12|\varphi''(x_0)|(x-x_0)^2\), so the
integrand is a Gaussian bell of width \(1/\sqrt{s|\varphi''(x_0)|}\). Replace \(g\) by \(g(x_0)\) and
integrate the Gaussian over the whole line:
\(\int e^{-\frac12s|\varphi''|(x-x_0)^2}dx=\sqrt{2\pi/(s|\varphi''|)}\).

\emph{Link to the general formula.} \eqref{eq:laplace-real} is \eqref{eq:sd-result} on the real line:
\(\varphi''(x_0)<0\) means \(\theta=\arg\varphi''(x_0)=\pi\), so \(\alpha=i\,e^{-i\pi/2}=i\cdot(-i)=1\). The real axis
is already the descent line. Chapter~2 uses \eqref{eq:laplace-real} for Stirling's formula.
\par\framebreak\par
\textbf{2. Stationary phase (an oscillating exponential).} Now the integrand is
\(g(x)\,e^{is\varphi(x)}\): it has size \(|g|\) and only spins around the unit circle.

\textbf{Intuition.} Where \(\varphi'\neq0\), the phase \(s\varphi\) changes quickly as \(x\) moves, so
\(e^{is\varphi}\) goes around the circle many times; neighbouring pieces point in opposite directions
and cancel. Only near a \emph{stationary point} \(x_0\), where \(\varphi'(x_0)=0\) and the phase is
momentarily flat, do the pieces line up and add. (This is the same reason light follows the
path of stationary optical length.)

\textbf{Result.} For a single stationary point \(x_0\) inside \((a,b)\), with \(\varphi''(x_0)\ne0\) and
\(g(x_0)\ne0\) (several stationary points add their contributions), and writing
\(\operatorname{sgn}\varphi''(x_0)=\pm1\) for the sign of \(\varphi''(x_0)\),
\begin{equation}
  \int_a^b g(x)\,e^{is\varphi(x)}\,dx
  \sim
  g(x_0)\,e^{is\varphi(x_0)}\,
  e^{i(\pi/4)\,\operatorname{sgn}\varphi''(x_0)}
  \sqrt{\frac{2\pi}{s\bigl|\varphi''(x_0)\bigr|}},
  \qquad s\to+\infty.
  \label{eq:stationary-phase}
\end{equation}
It looks like Laplace's formula, with one extra phase factor \(e^{\pm i\pi/4}\). The next frame
shows where that comes from.
\end{frame}

\begin{frame}[allowframebreaks=1.00]{Where the stationary-phase factor $e^{\pm i\pi/4}$ comes from}
\footnotesize

\textbf{Step 1: expand the phase near \(x_0\).} Since \(\varphi'(x_0)=0\),
\[
  \varphi(x)\approx\varphi(x_0)+\tfrac12\varphi''(x_0)\,(x-x_0)^2 .
\]
\textbf{Step 2: zoom in.} Only a window of width about \(1/\sqrt{s|\varphi''(x_0)|}\) contributes;
outside it the phase spins and cancels. So replace \(g\) by \(g(x_0)\), and rescale by
\(u=\sqrt s\,(x-x_0)\), \(dx=du/\sqrt s\). Then
\(s\cdot\tfrac12\varphi''(x_0)(x-x_0)^2=s\cdot\tfrac12\varphi''(x_0)\,u^2/s=\tfrac12\varphi''(x_0)\,u^2\). Letting the
\(u\)-limits run to \(\pm\infty\),
\[
  \int_a^b g(x)\,e^{is\varphi(x)}\,dx
  \sim
  e^{is\varphi(x_0)}\,\frac{g(x_0)}{\sqrt s}
  \int_{-\infty}^{\infty}
    e^{i\,\frac12\varphi''(x_0)\,u^2}\,du .
\]
\textbf{Step 3: the Fresnel integral.} We need \(\int_{-\infty}^{\infty}e^{\pm iau^2}du\) with
\(a=\frac12|\varphi''(x_0)|>0\) and the sign of \(\varphi''(x_0)\). Take the \(+\) sign. The trick is to rotate
the path of integration by \(45^\circ\), so that the oscillation turns into decay (as in Example~2 of the
gamma integrals in Chapter~2). Put \(u=e^{i\pi/4}v\) with \(v\) real:
\[
 u^2=e^{i\pi/2}v^2=i\,v^2,\qquad
 iau^2=ia\cdot iv^2=-a\,v^2,\qquad
 du=e^{i\pi/4}\,dv .
\]
\par\framebreak\par
\textbf{Step 3, continued.} So the oscillating integral becomes an ordinary Gaussian (using
\(\int e^{-av^2}dv=\sqrt{\pi/a}\)):
\[
 \int_{-\infty}^{\infty}e^{iau^2}\,du
 =e^{i\pi/4}\int_{-\infty}^{\infty}e^{-av^2}\,dv
 =e^{i\pi/4}\sqrt{\frac\pi a}.
\]
(Rotating the path is allowed by Cauchy's theorem: the arcs that close the sector contribute
nothing, as in Chapter~2.) For the \(-\) sign, rotate the other way, \(u=e^{-i\pi/4}v\), and get
\(e^{-i\pi/4}\sqrt{\pi/a}\). In one line:
\[
 \int_{-\infty}^{\infty}e^{\pm iau^2}\,du
 =\sqrt{\frac\pi a}\;e^{\pm i\pi/4},\qquad a>0 .
\]
\textbf{Step 4: assemble.} With \(a=\tfrac12|\varphi''(x_0)|\), \(\sqrt{\pi/a}=\sqrt{2\pi/|\varphi''(x_0)|}\), so
\[
 \frac{g(x_0)}{\sqrt s}\cdot\sqrt{\frac{2\pi}{|\varphi''(x_0)|}}\cdot e^{\pm i\pi/4}
 =g(x_0)\sqrt{\frac{2\pi}{s|\varphi''(x_0)|}}\;e^{i(\pi/4)\operatorname{sgn}\varphi''(x_0)} .
\]
Multiplying by \(e^{is\varphi(x_0)}\) gives \eqref{eq:stationary-phase}. The phase \(e^{\pm i\pi/4}\) is simply the
\(45^\circ\) tilt of the path that turns the oscillation into decay. Chapter~3 applies this rule to
the Airy integral at large negative argument.
\end{frame}

\begin{frame}[allowframebreaks=1.00]{Worked example: large binomial coefficient}
\footnotesize

\textbf{Goal.} Find how the binomial coefficient \(\binom nm=\frac{n!}{m!\,(n-m)!}\) behaves when \(n\) is
large and the ratio \(t=m/n\) is fixed (with \(0<t<1\)). To use steepest descent we must first write
\(\binom nm\) as an integral of the form \eqref{eq:sd-int}, with \(s=n\).

\textbf{Step 1: a contour integral for \(\binom nm\).} By the binomial theorem,
\((1+z)^n=\sum_{k=0}^n\binom nk z^k\). So \(\binom nm\) is the coefficient of \(z^m\), and the coefficient formula
\eqref{eq:laurent-coef} (center \(0\), any counterclockwise circle around \(0\)) gives
\[
  \binom nm=\frac1{2\pi i}\oint\frac{(1+z)^n}{z^{m+1}}\,dz .
\]
\textbf{Step 2: bring it to the form \(g\,e^{nf}\).} Split \(\frac{(1+z)^n}{z^{m+1}}=\frac1z\cdot(1+z)^n\cdot z^{-m}\) and
write the powers as exponentials, using \(m=nt\):
\[
 (1+z)^n=e^{n\log(1+z)},\qquad z^{-m}=e^{-m\log z}=e^{-nt\log z}.
\]
So
\[
  \binom nm=\oint g(z)\,e^{nf(z)}\,dz,
  \qquad
  g(z)=\frac{1}{2\pi i z},
  \qquad
  f(z)=\log(1+z)-t\log z .
\]
(The logarithms are taken on branches that are smooth near the positive real axis, where the
saddle will be.)
\par\framebreak\par
\textbf{Step 3: find the saddle.} Differentiate:
\[
  f'(z)=\frac1{1+z}-\frac t z,
  \qquad
  f''(z)=-\frac1{(1+z)^2}+\frac t{z^2}.
\]
Set \(f'(z_0)=0\) and multiply by \(z_0(1+z_0)\): \(z_0-t(1+z_0)=0\), so \(z_0(1-t)=t\):
\[
  z_0=\frac{t}{1-t}=\frac{m}{n-m}>0,\qquad 1+z_0=\frac{1}{1-t}.
\]
\textbf{Step 4: the height and curvature at the saddle.} Using \(\log(1+z_0)=-\log(1-t)\) and
\(\log z_0=\log t-\log(1-t)\):
\begin{align*}
  f(z_0)
  &=-\log(1-t)-t\bigl[\log t-\log(1-t)\bigr]
   =-t\log t-(1-t)\log(1-t),\\
  f''(z_0)
  &=-(1-t)^2+\frac{t\,(1-t)^2}{t^2}
   =(1-t)^2\Bigl(\frac1t-1\Bigr)
   =\frac{(1-t)^3}{t}>0.
\end{align*}
\par\framebreak\par
\textbf{Step 5: choose the path through the saddle.} \(f''(z_0)\) is positive real, so \(\theta=0\), and by
\eqref{eq:descent-direction} the descent direction is \(\phi=\pi/2\): vertical. The circle \(|z|=z_0\)
crosses the real axis vertically at \(z_0\), so it is a natural path. Two checks:
\begin{itemize}
  \item \emph{\(z_0\) is the highest point on this circle.} On \(z=z_0e^{i\psi}\), the size of the integrand
        is \(|1+z|^n/(2\pi z_0^{m+1})\), and \(|z|=z_0\) is fixed, so only \(|1+z|\) varies:
        \[
         |1+z_0e^{i\psi}|^2=(1+z_0\cos\psi)^2+(z_0\sin\psi)^2=1+2z_0\cos\psi+z_0^2 ,
        \]
        which is largest at \(\psi=0\), that is, at \(z=z_0\).
  \item \emph{Direction.} Going counterclockwise, the circle passes \(z_0\) moving upward, in the
        direction \(+i\). This agrees with \(\alpha=i\,e^{-i\theta/2}=i\), so the sign in \eqref{eq:sd-result} is
        right.
\end{itemize}
\par\framebreak\par
\textbf{Step 6: apply the formula \eqref{eq:sd-result} with \(s=n\).} With \(\theta=0\), \(e^{-i\theta/2}=1\), and
\(g(z_0)=1/(2\pi iz_0)\); the factor \(i\) cancels the \(i\) in \(g(z_0)\):
\[
  \binom nm
  \sim
  \frac{1}{2\pi i z_0}\,e^{nf(z_0)}\cdot i\cdot
  \sqrt{\frac{2\pi}{n f''(z_0)}}
  =
  \frac{1}{2\pi z_0}\,\sqrt{\frac{2\pi}{n f''(z_0)}}\;e^{nf(z_0)} .
\]
\textbf{Step 7: simplify the prefactor.} Insert \(1/z_0=(1-t)/t\) and \(1/f''(z_0)=t/(1-t)^3\), and bring
everything under one square root:
\[
  \frac{1-t}{2\pi t}
    \sqrt{\frac{2\pi t}{n(1-t)^3}}
  =\sqrt{\frac{(1-t)^2}{(2\pi t)^2}\cdot\frac{2\pi t}{n(1-t)^3}}
  =\sqrt{\frac{1}{2\pi\,n\,t\,(1-t)}} .
\]
\textbf{Result.} With \(e^{nf(z_0)}=\exp\bigl[-n\bigl(t\log t+(1-t)\log(1-t)\bigr)\bigr]\),
\begin{equation}
  \binom nm
  \sim
  \frac{\exp\bigl[-n\bigl(t\log t+(1-t)\log(1-t)\bigr)\bigr]}
       {\sqrt{2\pi n t(1-t)}},\qquad t=\frac mn .
  \label{eq:binom-sd}
\end{equation}
\par\framebreak\par
\textbf{Check: \(n=20\), \(m=10\)} (\(t=\tfrac12\)). The exponent is
\(-20\bigl(\tfrac12\log\tfrac12+\tfrac12\log\tfrac12\bigr)=20\log2\), so the numerator is \(2^{20}=1\,048\,576\);
the denominator is \(\sqrt{2\pi\cdot20\cdot\frac14}=\sqrt{10\pi}\approx5.605\). The formula gives
\(\approx187\,079\); the exact value is \(184\,756\), only \(1.3\%\) less. (Stirling's formula of Chapter~2,
applied to each factorial, gives the same result.)

\textbf{Meaning: entropy.} \(H(t)=-t\log t-(1-t)\log(1-t)\) is the \emph{entropy} of a coin with
heads-probability \(t\), and \eqref{eq:binom-sd} says \(\binom nm\approx e^{nH(t)}/\sqrt{2\pi nt(1-t)}\). The chance
of one particular sequence with \(m\) heads is
\(t^m(1-t)^{n-m}=e^{m\log t+(n-m)\log(1-t)}=e^{-nH(t)}\). Multiplying, the probability of exactly \(m\)
heads is
\[
 \binom nm t^m(1-t)^{n-m}\approx\frac{1}{\sqrt{2\pi nt(1-t)}} ,
\]
the height of the peak of a Gaussian with the binomial variance \(nt(1-t)\).
\end{frame}

\begin{frame}[allowframebreaks=1.00]{Summary}
\footnotesize

\begin{itemize}
 \item Complex differentiability requires one direction-independent difference-quotient
       limit.  With continuous first partial derivatives, this is equivalent to the
       Cauchy--Riemann equations $u_x=v_y$, $u_y=-v_x$.
 \item Cauchy's theorem: a holomorphic $f$ integrates to zero around a closed contour, so
       contours may be deformed freely through regions where $f$ is holomorphic.  Cauchy's
       integral formula yields all derivatives and the local Taylor expansion.
 \item A Laurent series is valid on an annulus.  Its principal part classifies an isolated
       singularity as removable, a pole, or essential.  Branch points are not isolated and
       do not possess ordinary residues.
\framebreak
 \item The residue is the coefficient $a_{-1}$.  The residue theorem evaluates a closed
       contour integral as $2\pi i$ times the sum of enclosed residues.
 \item Three real-integral templates: semicircle (rational/exponential), keyhole
       (powers and branch cuts, as in $\int_0^\infty x^a/(1+x)^2\,dx$), and Jordan's
       indented semicircle (oscillatory integrals such as $\int\sin x/x\,dx$).  Always
       check orientation and vanishing arcs before taking limits.
 \item \textbf{Steepest descent} \eqref{eq:sd-result} is a local Gaussian approximation at a
       nondegenerate dominant saddle.  Real special cases: Laplace method
       \eqref{eq:laplace-real} (Stirling in Ch.~2) and stationary phase
       \eqref{eq:stationary-phase} (Airy in Ch.~3); large-$x$ Bessel (Ch.~5) uses the complex form.  Orientation and the phase
       of the square root are part of the answer.
\end{itemize}
\end{frame}

\end{document}
